\section{2022 Mathematical Physics}

\textit{Solve every problem.}

\begin{exercise}
\label{physics:2022:1}
\begin{enumerate}[label=(\alph*)]
    \item A symmetry transformation in quantum mechanics is represented by a unitary or anti-unitary operator acting on a Hilbert space. The time reversal transformation $\Theta$ relates the wave function at time $t$ to time $-t$. Prove: $\Theta$ is an anti-unitary operator.
    \item Consider state vector $|\psi\rangle$ for a quantum system. A time reversal transformation is represented by an anti-unitary operator $\Theta$. We now consider position space wavefunction $\psi(x) = \langle x | \psi\rangle$, and $\Theta |x\rangle = |x\rangle$. Prove: the position space wave function for $\Theta|\psi\rangle$ is $\psi(x)^*$.\sidenote{The asterisk $^*$ denotes \textbf{complex conjugation}. For any complex number $z = a + bi$ (where $a, b \in \mathbb{R}$), its conjugate is $z^* = a - bi$. In quantum mechanics, wavefunctions are complex-valued, and the conjugate is essential for defining probabilities ($|\psi|^2 = \psi^* \psi$) and inner products ($\langle \phi | \psi \rangle^* = \langle \psi | \phi \rangle$).}
    \item A one-dimensional quantum system is invariant under time reversal transformation, and so its Hamiltonian\sidenote{$H$ is the \textbf{Hamiltonian operator}, representing the total energy of the system. The constant $\hbar = h/2\pi$ is the \textbf{reduced Planck constant}.} satisfies $\Theta H = H \Theta$. If an energy eigenstate $|\psi\rangle$ has no degeneracy, prove: it is possible to take the position space energy eigenfunction to be real: $\psi(x)^* = \psi(x)$.
\end{enumerate}
\end{exercise}

\begin{solution}
\textbf{Prerequisites / Core Concepts:}
\begin{itemize}
    \item \textbf{Anti-unitary Operator:} An operator $\Theta$ that is anti-linear ($\Theta(c_1|\psi\rangle + c_2|\phi\rangle) = c_1^*\Theta|\psi\rangle + c_2^*\Theta|\phi\rangle$) and satisfies the inner product property $\langle \Theta \psi | \Theta \phi \rangle = \langle \psi | \phi \rangle^* = \langle \phi | \psi \rangle$.
    \item \textbf{Eigenstate and Eigenvalue:} For a linear operator $A$, a non-zero state vector $|\psi\rangle$ is an \textbf{eigenstate} if $A|\psi\rangle = \lambda |\psi\rangle$ for some scalar $\lambda$, called the \textbf{eigenvalue}. Geometrically, the operator $A$ does not change the "direction" of $|\psi\rangle$ in Hilbert space, only scales it by $\lambda$.
    \item \textbf{Energy Eigenstate:} A specific type of eigenstate where the operator is the Hamiltonian $H$. It satisfies the \textbf{time-independent Schrödinger equation}: $H|\psi\rangle = E|\psi\rangle$, where $E$ is the energy eigenvalue.
    \item \textbf{Schrödinger Equation:} $\mathrm{i}\hbar \frac{\partial |\psi\rangle}{\partial t} = H|\psi\rangle$.\sidenote{The Hamiltonian $H$ defines the energy of the system, while $\hbar$ relates the time evolution to the energy scale.} This is the fundamental equation that governs the time evolution of quantum states.
    \item \textbf{Quantum Degeneracy:} A situation where multiple linearly independent eigenstates correspond to the same energy eigenvalue. \footnote{\textbf{Example of Degeneracy:} Consider a particle in a 2D square box of side length $L$. The energy levels are $E_{n_x, n_y} = \frac{\hbar^2 \pi^2}{2m L^2} (n_x^2 + n_y^2)$. The states $(n_x=1, n_y=2)$ and $(n_x=2, n_y=1)$ both have the same energy $E = \frac{5\hbar^2 \pi^2}{2m L^2}$, but their wavefunctions $\psi_{1,2}(x,y)$ and $\psi_{2,1}(x,y)$ are distinct. This is a 2-fold degeneracy arising from the spatial symmetry of the box.}
    \item \textbf{Non-degeneracy:} A condition where each energy eigenvalue $E$ corresponds to exactly one linearly independent eigenstate $|\psi\rangle$. Mathematically, this means the eigenspace $V_E = \{ |\phi\rangle : H|\phi\rangle = E|\phi\rangle \}$ is 1-dimensional.\sidenote{In simple matrix terms: the Hamiltonian is a matrix and the states are eigenvectors. \textbf{Non-degeneracy} means that for a specific eigenvalue (number), there is only \textbf{one} unique eigenvector (arrow). If you can find two different, independent arrows that are both scaled by the same number, that number is \textbf{degenerate}. Non-degenerate means the 'energy level' is not shared by multiple distinct states.} Any two states with the same energy $E$ must be proportional to each other ($|\phi_1\rangle = c |\phi_2\rangle$).
\end{itemize}

\textbf{Intuition / Motivation:}
In classical mechanics, if we reverse time ($t \to -t$), the position remains the same but the velocity (and momentum) reverses direction. In quantum mechanics, time evolution is governed by the imaginary unit $i$ in the Schrödinger equation. To reverse time and keep the physics consistent, we must also reverse the sign of $i$ (complex conjugation), which naturally makes the time reversal operator anti-unitary. For non-degenerate energy eigenstates, time reversal cannot change the state to a new independent state, so it can only change its phase. By choosing the overall phase carefully, we can make the wavefunction entirely real, which means it looks exactly like its own time-reversed version.

\textbf{Logical Step-by-Step Breakdown:}

\textbf{(a) Proof of Anti-unitarity:}
\begin{enumerate}
    \item \textbf{Step 1: Setup the time-reversed Schrödinger equation.}
    The system state $|\psi(t)\rangle$ satisfies the Schrödinger equation:\sidenote{In this equation, $H$ is the \textbf{Hamiltonian operator}, which represents the total energy of the quantum system (the sum of kinetic and potential energies). The constant $\hbar = h / 2\pi$ is the \textbf{reduced Planck constant}, which defines the scale of quantum actions and relates energy to frequency.}
    \[ \mathrm{i}\hbar \frac{\partial |\psi(t)\rangle}{\partial t} = H |\psi(t)\rangle \]
    Under time reversal $t \to -t$, let the time-reversed state be $|\psi'(t)\rangle = \Theta |\psi(-t)\rangle$. We want this new state to satisfy the same Schrödinger equation (assuming the Hamiltonian $H$ is invariant under time reversal):
    \[ \mathrm{i}\hbar \frac{\partial |\psi'(t)\rangle}{\partial t} = H |\psi'(t)\rangle \]
    
    \item \textbf{Step 2: Apply the time reversal operator to the original equation.}
    Starting with the original Schrödinger equation evaluated at time $-t$:
    \[ \mathrm{i}\hbar \frac{\partial |\psi(-t)\rangle}{\partial (-t)} = H |\psi(-t)\rangle \]
    This simplifies to:
    \[ -\mathrm{i}\hbar \frac{\partial |\psi(-t)\rangle}{\partial t} = H |\psi(-t)\rangle \]
    Now apply the time reversal operator $\Theta$ to both sides. Assuming $\Theta$ commutes with the Hamiltonian $H$ and the time derivative $\frac{\partial}{\partial t}$:
    \[ \Theta \left( -\mathrm{i}\hbar \frac{\partial |\psi(-t)\rangle}{\partial t} \right) = \Theta H |\psi(-t)\rangle = H \Theta |\psi(-t)\rangle = H |\psi'(t)\rangle \]
    
    \item \textbf{Step 3: Compare and deduce anti-linearity.}
    We require the time-reversed state to satisfy the Schrödinger equation: $\mathrm{i}\hbar \frac{\partial |\psi'(t)\rangle}{\partial t} = H |\psi'(t)\rangle$. Substituting this into the right-hand side of our derived equation gives:
    \[ \Theta \left( -\mathrm{i}\hbar \frac{\partial |\psi(-t)\rangle}{\partial t} \right) = \mathrm{i}\hbar \frac{\partial |\psi'(t)\rangle}{\partial t} = \mathrm{i}\hbar \frac{\partial \Theta |\psi(-t)\rangle}{\partial t} \]
    For this equality to hold for arbitrary states, the operator $\Theta$ must pull out the constant $-i$ as its complex conjugate $+i$. This means $\Theta(i) = -i \Theta$. Since it anti-commutes with $i$, $\Theta$ is an anti-linear operator. To preserve probabilities ($|\langle \Theta\psi | \Theta\phi \rangle| = |\langle \psi | \phi \rangle|$), it must also be anti-unitary.
\end{enumerate}

\textbf{(b) Wavefunction Transformation:}
\begin{enumerate}
    \item \textbf{Step 1: Define the target wavefunction.}
    Let $|\psi'\rangle = \Theta |\psi\rangle$ be the time-reversed state. We want to find its position space representation, which is by definition:
    \[ \psi'(x) = \langle x | \psi'\rangle = \langle x | \Theta \psi \rangle \]

    \item \textbf{Step 2: Relate bras to kets under anti-unitary transformation.}
    Recall the fundamental definition of an anti-unitary operator $\Theta$: for any two states $|\phi\rangle$ and $|\xi\rangle$, the inner product of their transformed versions is the complex conjugate of their original inner product:
    \[ \langle \Theta \phi | \Theta \xi \rangle = \langle \phi | \xi \rangle^* = \langle \xi | \phi \rangle \]
    This is the \textbf{anti-unitary identity}.\sidenote{For a normal unitary operator $U$, we have $\langle U\phi | U\xi \rangle = \langle \phi | \xi \rangle$. The "swapping" to $\langle \xi | \phi \rangle$ for $\Theta$ is a direct consequence of its anti-linearity. It ensures that the physics (probabilities) remains invariant even though the "arrows" of the complex amplitudes are reversed.} It allows us to move $\Theta$ from the ket to the bra, provided we "reverse" the order of the inner product.

    \item \textbf{Step 3: Utilize the given condition $\Theta |x\rangle = |x\rangle$.}
    We are told that the position basis states $|x\rangle$ are invariant under time reversal. In Dirac notation, this means the ket $|x\rangle$ is the same as the ket $\Theta |x\rangle$. Consequently, the corresponding bra $\langle x |$ is the same as the bra $\langle \Theta x |$.\sidenote{The symbol $\langle \Theta x |$ represents the \textbf{dual vector} (bra) corresponding to the transformed state $\Theta |x\rangle$. In Hilbert space theory, if $\Theta$ were a linear operator $A$, then $\langle A x |$ would be the bra $\langle x | A^\dagger$. However, since $\Theta$ is \textbf{anti-linear}, the relationship is defined by the property $\langle \Theta x | \Theta \psi \rangle = \langle \psi | x \rangle$, which ensures that the inner product (a complex amplitude) is properly conjugated to maintain physical consistency.}
    Substituting this into our target expression:
    \[ \psi'(x) = \langle \Theta x | \Theta \psi \rangle \]

    \item \textbf{Step 4: Apply the anti-unitary identity.}
    Now we apply the identity from Step 2, using $|\phi\rangle = |x\rangle$ and $|\xi\rangle = |\psi\rangle$:
    \[ \langle \Theta x | \Theta \psi \rangle = \langle \psi | x \rangle \]

    \item \textbf{Step 5: Final conclusion.}
    By the standard definition of the wavefunction, $\langle x | \psi \rangle = \psi(x)$. Therefore, $\langle \psi | x \rangle = \langle x | \psi \rangle^* = \psi(x)^*$.
    Combining everything, we reach the final result:
    \[ \psi'(x) = \psi(x)^* \]
    This means that time reversal effectively replaces the original wavefunction with its complex conjugate at every point in space.
\end{enumerate}
\textbf{(c) Reality of Eigenfunctions:}
\begin{enumerate}
    \item \textbf{Step 1: Show $\Theta|\psi\rangle$ is an energy eigenstate with the same energy.}
    Given the Hamiltonian commutes with time reversal, $[H, \Theta] = 0$. Let $|\psi\rangle$ be an energy eigenstate with eigenvalue $E$, so $H|\psi\rangle = E|\psi\rangle$.
    Apply $H$ to $\Theta|\psi\rangle$:
    \[ H(\Theta|\psi\rangle) = \Theta(H|\psi\rangle) = \Theta(E|\psi\rangle) \]
    Since $\Theta$ is anti-linear and $E$ is real, $\Theta(E|\psi\rangle) = E(\Theta|\psi\rangle)$.
    Thus, $\Theta|\psi\rangle$ is also an eigenstate with energy $E$.
    
    \item \textbf{Step 2: Utilize the non-degeneracy condition.}
    The problem states that the state $|\psi\rangle$ has no degeneracy. This means any state with energy $E$ must be a scalar multiple of $|\psi\rangle$. Therefore:
    \[ \Theta|\psi\rangle = c|\psi\rangle \]
    Since $\Theta$ preserves the norm, $|c|^2 = 1$, so we can write $c = e^{\mathrm{i} \alpha}$ for some real phase factor $\alpha$:
    \[ \Theta|\psi\rangle = e^{\mathrm{i} \alpha} |\psi\rangle \]
    
    \item \textbf{Step 3: Relate this to the position space wavefunction.}
    Taking the inner product with $\langle x |$ on both sides:
    \[ \langle x | \Theta | \psi \rangle = \langle x | e^{\mathrm{i} \alpha} | \psi \rangle \]
    \sidenote{The formal mathematical definition of the matrix element $\langle \phi | \hat{A} | \psi \rangle$ is the \textbf{inner product} $(\phi, \hat{A}\psi)$ in the Hilbert space $\mathcal{H}$. Specifically, $\langle x | \Theta | \psi \rangle \equiv (x, \Theta \psi)$. This represents the value of the linear functional $\langle x |$ evaluated at the vector $\Theta |\psi\rangle$, which is by definition the position-space representation (wavefunction) of the state $\Theta |\psi\rangle$ at point $x$.}
    From part (b), the left side is $\psi(x)^*$. The right side is $e^{\mathrm{i} \alpha} \langle x | \psi \rangle = e^{\mathrm{i} \alpha} \psi(x)$. So:
    \[ \psi(x)^* = e^{\mathrm{i} \alpha} \psi(x) \]
    
    \item \textbf{Step 4: Construct a real eigenfunction.}
    We can define a new wavefunction by multiplying $\psi(x)$ by a constant phase factor to absorb the $e^{\mathrm{i} \alpha}$. Let $\phi(x) = e^{\mathrm{i} \alpha/2} \psi(x)$.
    Now, check the complex conjugate of $\phi(x)$:
    \[ \phi(x)^* = (e^{\mathrm{i} \alpha/2} \psi(x))^* = e^{-i\alpha/2} \psi(x)^* \]
    Substitute $\psi(x)^* = e^{\mathrm{i} \alpha} \psi(x)$:
    \[ \phi(x)^* = e^{-i\alpha/2} (e^{\mathrm{i} \alpha} \psi(x)) = e^{\mathrm{i} \alpha/2} \psi(x) = \phi(x) \]
    Because $\phi(x)^* = \phi(x)$, the redefined energy eigenfunction $\phi(x)$ is purely real.
\end{enumerate}
\end{solution}

\begin{exercise}
\label{physics:2022:2}
Consider following quantum Hamiltonian:\sidenote{Here, $\omega$ is the \textbf{angular frequency} of the harmonic oscillator. In classical mechanics, it is related to the spring constant $k$ and mass $m$ by $\omega = \sqrt{k/m}$. In quantum mechanics, it determines the spacing between energy levels, $\Delta E = \hbar\omega$.}
\begin{equation}
    H_0 = \frac{p_1^2}{2m} + \frac{1}{2}m\omega^2 x_1^2 + \frac{p_2^2}{2m} + \frac{1}{2}m\omega^2 x_2^2
\end{equation}
This is the Hamiltonian for two decoupled harmonic oscillators.

\begin{enumerate}[label=(\alph*)]
    \item Calculate the eigenstates and eigenvalues for $H_0$ (an energy eigenstate could be labeled as $|n_1, n_2\rangle$).
    \item Assume the creation and annihilation operators for two harmonic oscillators are $a_i^\dagger, a_i, i=1,2$.\sidenote{The \textbf{annihilation operator} $a_i$ and \textbf{creation operator} $a_i^\dagger$ (also called ladder operators) are defined in terms of position and momentum as $a_i = \sqrt{\frac{m\omega}{2\hbar}}(x_i + \frac{\mathrm{i}}{m\omega}p_i)$ and $a_i^\dagger = \sqrt{\frac{m\omega}{2\hbar}}(x_i - \frac{\mathrm{i}}{m\omega}p_i)$. They satisfy the \textbf{canonical commutation relation} $[a_i, a_j^\dagger] = \delta_{ij}$. Physically, $a_i$ "lowers" the energy level of the $i$-th oscillator by $\hbar\omega$, while $a_i^\dagger$ "raises" it.} Define following operators
    \begin{equation}
        J_+ = a_1^\dagger a_2, \quad J_- = a_2^\dagger a_1, \quad J_z = \frac{1}{2}(a_1^\dagger a_1 - a_2^\dagger a_2)
    \end{equation}
    \begin{enumerate}[label=\roman*.]
        \item Prove that: $[J_z, J_\pm] = \pm J_\pm$, $[J_+, J_-] = 2J_z$.
        \item Consider one eigenvalue $E_n$ of $H_0$, (here $n_1 + n_2 = n$). Prove that: all eigenstates of $E_n$ form an irreducible representation of $su(2)$ Lie algebra, and compute the spin.\sidenote{In the \textbf{Schwinger boson representation}, a system of two oscillators can represent a "spin" system. The \textbf{spin} (total angular momentum quantum number $j$) is defined by the total number of quanta $n$ through the relation $n = 2j$, or $j = n/2$. This means a system with $n$ total energy units behaves mathematically like a particle with spin $j=n/2$.}
    \end{enumerate}
    \item Consider following perturbed Hamiltonian ($\lambda$ is small)
    \begin{equation}
        H = H_0 + \lambda x_1^2 p_2^2
    \end{equation}
    Compute the first order correction to the energy for the energy level $n_1 + n_2 = 2$.
\end{enumerate}
\end{exercise}

\begin{solution}
\textbf{Intuition / Motivation:}
This problem explores the deep connection between the 2D harmonic oscillator and the theory of angular momentum ($su(2)$ symmetry). We start with two independent (decoupled) oscillators. Because they have the same frequency $\omega$, the system possesses a hidden symmetry: we can "rotate" energy between the two oscillators without changing the total energy. This symmetry is mathematically captured by operators $J_\pm$ and $J_z$, which satisfy the same rules as physical spin or orbital angular momentum. In the final part, we introduce a perturbation that breaks this symmetry and mixes the states, requiring us to use degenerate perturbation theory to find how the energy levels split.

\textbf{Logical Step-by-Step Breakdown:}

\textbf{(a) Eigenstates and Eigenvalues of $H_0$}
The unperturbed Hamiltonian $H_0$ describes a 2D isotropic harmonic oscillator. It is the sum of two independent 1D Hamiltonians:
\[ H_0 = \underbrace{\left( \frac{p_1^2}{2m} + \frac{1}{2}m\omega^2 x_1^2 \right)}_{H_1} + \underbrace{\left( \frac{p_2^2}{2m} + \frac{1}{2}m\omega^2 x_2^2 \right)}_{H_2} = H_1 + H_2. \]
Here, $H_1$ and $H_2$ are \textbf{operators}, specifically the energy operators (Hamiltonians) for the $x_1$ and $x_2$ degrees of freedom. 

\textbf{Separation of Variables:}
Because $H_1$ depends only on $(x_1, p_1)$ and $H_2$ depends only on $(x_2, p_2)$, they act on independent parts of the Hilbert space. This implies they \textbf{commute}: $[H_1, H_2] = 0$. In quantum mechanics, commuting operators share a common set of eigenfunctions. Therefore, we can solve the Schrödinger equation $H_0 \Psi = E \Psi$ by \textbf{separation of variables}:
\begin{enumerate}
    \item We assume the total wavefunction is a product: $\Psi(x_1, x_2) = \psi_1(x_1) \psi_2(x_2)$.
    \item Substituting into $(H_1 + H_2) \psi_1 \psi_2 = E \psi_1 \psi_2$ and dividing by $\psi_1 \psi_2$ gives:
    \[ \frac{H_1 \psi_1(x_1)}{\psi_1(x_1)} + \frac{H_2 \psi_2(x_2)}{\psi_2(x_2)} = E. \]
    \item Since the first term depends only on $x_1$ and the second only on $x_2$, each must be a constant: $E_1$ and $E_2$, where $E_1 + E_2 = E$.
\end{enumerate}
This reduces the 2D problem into two identical 1D harmonic oscillator problems.

\textbf{1. Derivation of the 1D Oscillator Spectrum:}
To find the energy eigenvalues of a single oscillator $H_i = \frac{p_i^2}{2m} + \frac{1}{2}m\omega^2 x_i^2$, we use the \textbf{algebraic ladder operator method}. We define the \textbf{annihilation} ($a_i$) and \textbf{creation} ($a_i^\dagger$) operators as follows:
\[ a_i = \sqrt{\frac{m\omega}{2\hbar}} \left( x_i + \frac{i}{m\omega} p_i \right), \quad a_i^\dagger = \sqrt{\frac{m\omega}{2\hbar}} \left( x_i - \frac{i}{m\omega} p_i \right). \]
Using the canonical commutation relation $[x_i, p_j] = i\hbar \delta_{ij}$, we compute the commutator of these ladder operators:
\[ [a_i, a_j^\dagger] = \frac{m\omega}{2\hbar} [x_i + \frac{i}{m\omega} p_i, x_j - \frac{i}{m\omega} p_j] = \frac{m\omega}{2\hbar} \delta_{ij} \left( -\frac{i}{m\omega}[x_i, p_i] + \frac{i}{m\omega}[p_i, x_i] \right). \]
Since $[x_i, p_i] = i\hbar$ and $[p_i, x_i] = -i\hbar$, the term in the parenthesis becomes $-i/m\omega(i\hbar) + i/m\omega(-i\hbar) = 2\hbar/m\omega$. Thus:
\[ [a_i, a_j^\dagger] = \frac{m\omega}{2\hbar} \delta_{ij} \left( \frac{2\hbar}{m\omega} \right) = \delta_{ij}. \]
\textbf{Detailed Derivation of the Hamiltonian Identity:}
We want to show why $H_i = \hbar\omega (N_i + 1/2)$. We start with the product $a_i^\dagger a_i$:
\[ \begin{aligned}
a_i^\dagger a_i &= \left[ \sqrt{\frac{m\omega}{2\hbar}} \left( x_i - \frac{i}{m\omega} p_i \right) \right] \left[ \sqrt{\frac{m\omega}{2\hbar}} \left( x_i + \frac{i}{m\omega} p_i \right) \right] \\
&= \frac{m\omega}{2\hbar} \left( x_i - \frac{i}{m\omega} p_i \right) \left( x_i + \frac{i}{m\omega} p_i \right)
\end{aligned} \]
Distributing the operators (carefully keeping their order, as $[x, p] \neq 0$):
\[ \begin{aligned}
a_i^\dagger a_i &= \frac{m\omega}{2\hbar} \left[ x_i^2 + x_i \left( \frac{i}{m\omega} p_i \right) - \left( \frac{i}{m\omega} p_i \right) x_i - \left( \frac{i}{m\omega} p_i \right) \left( \frac{i}{m\omega} p_i \right) \right] \\
&= \frac{m\omega}{2\hbar} \left[ x_i^2 + \frac{i}{m\omega} (x_i p_i - p_i x_i) - \frac{i^2}{m^2\omega^2} p_i^2 \right] \\
&= \frac{m\omega}{2\hbar} \left[ x_i^2 + \frac{i}{m\omega} [x_i, p_i] + \frac{1}{m^2\omega^2} p_i^2 \right]
\end{aligned} \]
Now, substitute the canonical commutator $[x_i, p_i] = i\hbar$:
\[ \begin{aligned}
a_i^\dagger a_i &= \frac{m\omega}{2\hbar} \left[ x_i^2 + \frac{i}{m\omega} (i\hbar) + \frac{1}{m^2\omega^2} p_i^2 \right] \\
&= \frac{m\omega}{2\hbar} \left[ x_i^2 - \frac{\hbar}{m\omega} + \frac{p_i^2}{m^2\omega^2} \right] \\
&= \frac{m\omega}{2\hbar} x_i^2 + \frac{p_i^2}{2m\hbar\omega} - \frac{1}{2}
\end{aligned} \]
To reconstruct the Hamiltonian $H_i = \frac{p_i^2}{2m} + \frac{1}{2}m\omega^2 x_i^2$, we multiply the expression by $\hbar\omega$:
\[ \hbar\omega (a_i^\dagger a_i) = \frac{1}{2}m\omega^2 x_i^2 + \frac{p_i^2}{2m} - \frac{1}{2}\hbar\omega = H_i - \frac{1}{2}\hbar\omega. \]
Rearranging for $H_i$, and noting that $N_i = a_i^\dagger a_i$ is the number operator:
\[ H_i = \hbar\omega \left( a_i^\dagger a_i + \frac{1}{2} \right) = \hbar\omega \left( N_i + \frac{1}{2} \right). \]
This constant $1/2$ shift is a direct result of the non-vanishing commutator $[x, p]$, representing the \textbf{zero-point energy} of the quantum oscillator.

\textbf{Why are the eigenvalues of $N_i$ non-negative integers?}
By definition, let $|n_i\rangle$ be an eigenstate of the number operator $N_i$ such that $N_i |n_i\rangle = n_i |n_i\rangle$. We can prove that $n_i$ must be a non-negative integer:
\begin{enumerate}
    \item \textbf{Non-negativity}: The norm squared of the state $a_i |n_i\rangle$ is:
    \[ || a_i |n_i\rangle ||^2 = \langle n_i | a_i^\dagger a_i | n_i \rangle = \langle n_i | N_i | n_i \rangle = n_i. \]
    In a physical Hilbert space, the norm squared must be non-negative, which implies $n_i \geq 0$.
    \item \textbf{Shifting property}: Using the commutator $[N_i, a_i] = -a_i$, we see that:
    \[ N_i (a_i |n_i\rangle) = (a_i N_i + [N_i, a_i]) |n_i\rangle = (a_i n_i - a_i) |n_i\rangle = (n_i - 1) (a_i |n_i\rangle). \]
    This shows that $a_i |n_i\rangle$ is an eigenstate of $N_i$ with eigenvalue $n_i - 1$.
    \item \textbf{Quantization}: If $n_i$ were not an integer (e.g., $n_i = 1.5$), repeatedly applying the annihilation operator $a_i$ would eventually produce a state with a negative eigenvalue ($0.5 \to -0.5$), which violates the norm condition in Step 1. The only way to prevent this is if the sequence terminates at $n_i = 0$, such that $a_i |0\rangle = 0$.
\end{enumerate}
Thus, the possible values for $n_i$ are the integers $\{0, 1, 2, \dots\}$.

\textbf{Final Step: From Operators to Energy Values}
To find the energy eigenvalue $E_{n_i}$, we apply the Hamiltonian $H_i$ to an eigenstate $|n_i\rangle$:
\[ \begin{aligned}
H_i |n_i\rangle &= \hbar\omega \left( N_i + \frac{1}{2} \right) |n_i\rangle \\
&= \hbar\omega \left( N_i |n_i\rangle + \frac{1}{2} |n_i\rangle \right) \\
&= \hbar\omega \left( n_i |n_i\rangle + \frac{1}{2} |n_i\rangle \right) \quad (\text{since } N_i |n_i\rangle = n_i |n_i\rangle) \\
&= \hbar\omega \left( n_i + \frac{1}{2} \right) |n_i\rangle.
\end{aligned} \]
Comparing this to the \textbf{Schrödinger eigenvalue equation} $H_i |n_i\rangle = E_{n_i} |n_i\rangle$, we conclude that the energy eigenvalues are:
\[ E_{n_i} = \hbar\omega \left( n_i + \frac{1}{2} \right). \]

\textbf{2. Total Eigenstates:}
The eigenstates of the total system are the tensor product of the individual states:
\[ |n_1, n_2\rangle = |n_1\rangle \otimes |n_2\rangle. \]
These states are uniquely identified by the number of energy quanta $n_1$ and $n_2$ in each oscillator.

\textbf{3. Total Eigenvalues:}
The total energy is the sum of the energies of the two oscillators:
\[ E_{n_1, n_2} = E_{n_1} + E_{n_2} = \hbar\omega(n_1 + 1/2) + \hbar\omega(n_2 + 1/2) = \hbar\omega(n_1 + n_2 + 1). \]
If we define $n = n_1 + n_2$ as the total number of quanta, then $E_n = \hbar\omega(n+1)$.

\textbf{(b) $su(2)$ Lie Algebra and Representations}

\textbf{i. Proving the Commutation Relations:}
We use the canonical commutation relations $[a_i, a_j^\dagger] = \delta_{ij}$ and $[a_i, a_j] = [a_i^\dagger, a_j^\dagger] = 0$.

\textbf{1. Commutator $[J_z, J_+]$:}
Recall $J_z = \frac{1}{2}(a_1^\dagger a_1 - a_2^\dagger a_2) = \frac{1}{2}(N_1 - N_2)$ and $J_+ = a_1^\dagger a_2$.
\[ [J_z, J_+] = \frac{1}{2} [N_1 - N_2, a_1^\dagger a_2] = \frac{1}{2} \left( [N_1, a_1^\dagger a_2] - [N_2, a_1^\dagger a_2] \right). \]
Using the identities $[AB, C] = A[B, C] + [A, C]B$ and $[A, BC] = [A, B]C + B[A, C]$, along with the basic property $[N_i, a_j^\dagger] = \delta_{ij} a_i^\dagger$, $[N_i, a_j] = -\delta_{ij} a_i$:
\begin{itemize}
    \item \textbf{Derivation of $[N_1, a_1^\dagger]$:}
    \[ [N_1, a_1^\dagger] = [a_1^\dagger a_1, a_1^\dagger] = a_1^\dagger [a_1, a_1^\dagger] + [a_1^\dagger, a_1^\dagger] a_1 = a_1^\dagger(1) + 0 = a_1^\dagger. \]
    Thus $[N_1, a_1^\dagger a_2] = [N_1, a_1^\dagger] a_2 = a_1^\dagger a_2$.
    
    \item \textbf{Derivation of $[N_2, a_2]$:}
    \[ [N_2, a_2] = [a_2^\dagger a_2, a_2] = a_2^\dagger [a_2, a_2] + [a_2^\dagger, a_2] a_2 = 0 + (-1)a_2 = -a_2. \]
    Thus $[N_2, a_1^\dagger a_2] = a_1^\dagger [N_2, a_2] = a_1^\dagger (-a_2) = -a_1^\dagger a_2$.
\end{itemize}
Thus, $[J_z, J_+] = \frac{1}{2} (a_1^\dagger a_2 - (-a_1^\dagger a_2)) = a_1^\dagger a_2 = J_+$.
By Hermitian conjugation, $[J_z, J_-] = -J_-$.

\textbf{2. Commutator $[J_+, J_-]$:}
$J_+ = a_1^\dagger a_2$ and $J_- = a_2^\dagger a_1$.
\[ [J_+, J_-] = [a_1^\dagger a_2, a_2^\dagger a_1] = a_1^\dagger [a_2, a_2^\dagger a_1] + [a_1^\dagger, a_2^\dagger a_1] a_2. \]
Expanding the nested commutators:
\begin{itemize}
    \item $a_1^\dagger [a_2, a_2^\dagger a_1] = a_1^\dagger [a_2, a_2^\dagger] a_1 = a_1^\dagger (1) a_1 = N_1$.
    \item $[a_1^\dagger, a_2^\dagger a_1] a_2 = a_2^\dagger [a_1^\dagger, a_1] a_2 = a_2^\dagger (-1) a_2 = -N_2$.
\end{itemize}
Thus, $[J_+, J_-] = N_1 - N_2 = 2 J_z$.
These three operators satisfy the standard commutation relations for the $su(2)$ Lie algebra.
\textbf{ii. Proof of Irreducible Representation and Computation of Spin:}

For a fixed energy $E_n$, the total number of quanta $n = n_1 + n_2$ is constant.
The possible states are $\{|n, 0\rangle, |n-1, 1\rangle, \dots, |0, n\rangle\}$.

\begin{definition}
A \textbf{Lie algebra} $\mathfrak{g}$ is a vector space over a field (typically $\mathbb{R}$ or $\mathbb{C}$) equipped with a binary operation $[\cdot, \cdot]: \mathfrak{g} \times \mathfrak{g} \to \mathfrak{g}$, called the \textbf{Lie bracket}, that satisfies:
\begin{itemize}
    \item \textbf{Bilinearity:} $[ax + by, z] = a[x, z] + b[y, z]$.
    \item \textbf{Alternating (Antisymmetry):} $[x, x] = 0$, which implies $[x, y] = -[y, x]$.
    \item \textbf{Jacobi Identity:} $[x, [y, z]] + [y, [z, x]] + [z, [x, y]] = 0$.
\end{itemize}
In Exercise 2, we are working with the $su(2)$ Lie algebra, defined by the relations seen in part (i).
\end{definition}

\begin{definition}
A \textbf{representation} of a Lie algebra $\mathfrak{g}$ on a vector space $V$ is a Lie algebra homomorphism $\pi: \mathfrak{g} \to \mathfrak{gl}(V)$. This means $\pi$ is a linear map that preserves the commutation relations:
...
\[ \pi([X, Y]) = [\pi(X), \pi(Y)] = \pi(X)\pi(Y) - \pi(Y)\pi(X) \]
for all $X, Y \in \mathfrak{g}$. In Exercise 2, the map $\pi$ takes the abstract $su(2)$ generators and assigns them to specific operators (matrices) $J_\pm, J_z$ acting on the Hilbert space of the oscillators.
\end{definition}

\begin{definition}
In the context of the $su(2)$ Lie algebra, the \textbf{spin $j$} is a non-negative integer or half-integer ($j \in \{0, 1/2, 1, 3/2, \dots\}$) that labels the irreducible representation of the algebra. It physically represents the \textbf{total angular momentum} (or pseudo-angular momentum in this oscillator mapping) of the system, such that the eigenvalue of the Casimir operator $\mathbf{J}^2 = J_x^2 + J_y^2 + J_z^2$ is fixed at $j(j+1)\hbar^2$.
\end{definition}

\begin{definition}
An \textbf{irreducible representation of spin $j$} of the $su(2)$ Lie algebra is a $(2j+1)$-dimensional vector space $V_j$ spanned by a basis $\{|j, m\rangle\}$ where $m \in \{-j, -j+1, \dots, j\}$. These states are simultaneous eigenstates of the Casimir operator $\mathbf{J}^2$ and the generator $J_z$:
\begin{itemize}
    \item $\mathbf{J}^2 |j, m\rangle = j(j+1) |j, m\rangle$
    \item $J_z |j, m\rangle = m |j, m\rangle$
\end{itemize}
The representation is \textbf{irreducible} because there is no non-trivial subspace of $V_j$ that remains invariant under the action of the operators $J_+, J_-$, and $J_z$.
\end{definition}

\begin{remark}
\textbf{The Origin of Half-Integer Spin:}
A common point of confusion is why $j$ can be a half-integer (like $1/2$) while orbital angular momentum $l$ is always an integer. 
\begin{itemize}
    \item \textbf{Algebraic Requirement}: The $su(2)$ commutation relations $[J_+, J_-] = 2J_z$ and $[J_z, J_\pm] = \pm J_\pm$ only require that the difference between the maximum and minimum weights ($j - (-j) = 2j$) must be a non-negative \textbf{integer}. This allows $j$ to be $0, 1/2, 1, 3/2, \dots$.
    \item \textbf{Oscillator Mapping}: In our 2D oscillator, $j$ is defined by the total number of quanta: $j = n/2 = (n_1 + n_2)/2$. 
    \begin{itemize}
        \item If the total quanta $n$ is \textbf{even}, $j$ is an \textbf{integer} (e.g., $n=2 \implies j=1$).
        \item If the total quanta $n$ is \textbf{odd}, $j$ is a \textbf{half-integer}. Specifically, for the first excited state $n=1$, we must have $j=1/2$ because the system has exactly one quantum of energy.
    \end{itemize}
    \item \textbf{Weight States}: For $n=1$, the states are $|1, 0\rangle$ and $|0, 1\rangle$. The $J_z$ eigenvalues are $m = (n_1 - n_2)/2$, giving $m = \pm 1/2$. This pair forms the weight system for the \textbf{fundamental representation} of $su(2)$, which describes physical Spin-1/2 particles (like electrons).
\end{itemize}

\textbf{Visualizing Spin and Weights:}
To make these abstract definitions concrete, we can map the total oscillator quanta $n = n_1 + n_2$ to the spin $j = n/2$. The $J_z$ eigenvalue (the \textbf{weight}) is $m = (n_1 - n_2)/2$.

\begin{enumerate}
    \item \textbf{Case $j=0$ (Ground state, $n=0$):}
    \begin{itemize}
        \item State: $|0, 0\rangle \to |j=0, m=0\rangle$.
        \item Diagram: $\bullet$ (Singlet).
    \end{itemize}
    
    \item \textbf{Case $j=1/2$ (First excited state, $n=1$):}
    \begin{itemize}
        \item States: 
        $|1, 0\rangle \to |1/2, +1/2\rangle$ (Spin Up), 
        $|0, 1\rangle \to |1/2, -1/2\rangle$ (Spin Down).
        \item Weight Diagram: $\circ \xleftrightarrow{J_\pm} \circ$ (Doublet).
    \end{itemize}
    
    \item \textbf{Case $j=1$ (Second excited state, $n=2$):}
    \begin{itemize}
        \item States: 
        $|2, 0\rangle \to |1, +1\rangle$, 
        $|1, 1\rangle \to |1, 0\rangle$, 
        $|0, 2\rangle \to |1, -1\rangle$.
        \item Weight Diagram: $\circ \xleftrightarrow{J_\pm} \circ \xleftrightarrow{J_\pm} \circ$ (Triplet).
    \end{itemize}
\end{enumerate}
The raising operator $J_+ = a_1^\dagger a_2$ moves a quantum from oscillator 2 to oscillator 1, increasing $m$ by 1, while the lowering operator $J_- = a_2^\dagger a_1$ does the reverse.
\end{remark}

\textbf{Step-by-Step Verification of Irreducibility:}

We provide a formal verification that the energy eigenstates of $H_0$ for a fixed total energy $E_n$ form an irreducible representation of the $su(2)$ Lie algebra.

\begin{enumerate}
    \item \textbf{Step 1: Identify the Basis States.}
    As identified above, for a fixed energy $E_n = \hbar\omega(n+1)$, the possible states are:
    \[ \mathcal{B}_n = \{ |n, 0\rangle, |n-1, 1\rangle, \dots, |0, n\rangle \} \]
    The dimension of the subspace is $d = n+1$.

    \item \textbf{Step 2: Verify the Action of the su(2) Generators.}
    The generators $J_z, J_+, J_-$ act on these states as follows:
    \begin{itemize}
        \item $J_z |n_1, n_2\rangle = \frac{1}{2}(n_1 - n_2) |n_1, n_2\rangle$. These are the \textbf{weights} $m$.
        \item $J_+ |n_1, n_2\rangle = a_1^\dagger a_2 |n_1, n_2\rangle \propto |n_1+1, n_2-1\rangle$. This raises $m$ by 1.
        \item $J_- |n_1, n_2\rangle = a_2^\dagger a_1 |n_1, n_2\rangle \propto |n_1-1, n_2+1\rangle$. This lowers $m$ by 1.
    \end{itemize}
    Crucially, these operators preserve the total quanta $n = n_1 + n_2$, so they map the subspace into itself.

    \item \textbf{Step 3: Prove Irreducibility.}
    A representation is \textbf{irreducible} if every state in the subspace can be reached from any other state using the ladder operators $J_\pm$. 
    \begin{itemize}
        \item By applying $J_+ = a_1^\dagger a_2$ multiple times, we can move all energy quanta from the second oscillator to the first (reaching $|n, 0\rangle$).
        \item By applying $J_- = a_2^\dagger a_1$ multiple times, we can move all energy quanta from the first oscillator to the second (reaching $|0, n\rangle$).
    \end{itemize}
    Since every state $|n_1, n_2\rangle$ is connected in this chain, the representation is irreducible.

    \item \textbf{Step 4: Compute the Spin $j$.}
    In $su(2)$ theory, an irreducible representation of spin $j$ has a dimension of $2j + 1$. Equating this to our subspace dimension $d = n+1$:
    \[ 2j + 1 = n + 1 \implies 2j = n \implies \mathbf{j = \frac{n}{2}} \]
    This confirms that the eigenstates for total quanta $n$ form a \textbf{spin-$n/2$} representation.
\end{enumerate}
\textbf{(c) First-Order Correction for $n_1 + n_2 = 2$}
We use \textbf{degenerate perturbation theory} for the $n=2$ subspace. The energy level $E_2 = 3\hbar\omega$ is 3-fold degenerate.

\textbf{1. The Degenerate Basis:}
The states with total quanta $n=2$ are:
\[ \psi_1 = |2, 0\rangle, \quad \psi_2 = |1, 1\rangle, \quad \psi_3 = |0, 2\rangle. \]

\textbf{2. Step-by-Step Operator Expansion:}
The perturbation is $V = \lambda x_1^2 p_2^2$. We first express $x_1^2$ and $p_2^2$ in terms of ladder operators.
Using $x_1 = \sqrt{\frac{\hbar}{2m\omega}}(a_1 + a_1^\dagger)$ and $p_2 = \mathrm{i}\sqrt{\frac{m\hbar\omega}{2}}(a_2^\dagger - a_2)$:
\begin{itemize}
    \item \textbf{Expanding $x_1^2$:}
    \[ x_1^2 = \frac{\hbar}{2m\omega} (a_1 + a_1^\dagger)^2 = \frac{\hbar}{2m\omega} (a_1^2 + (a_1^\dagger)^2 + a_1 a_1^\dagger + a_1^\dagger a_1). \]
    Since $a_1 a_1^\dagger = a_1^\dagger a_1 + 1 = N_1 + 1$, we have:
    \[ x_1^2 = \frac{\hbar}{2m\omega} (a_1^2 + (a_1^\dagger)^2 + 2N_1 + 1). \]
    \item \textbf{Expanding $p_2^2$:}
    \[ p_2^2 = -\frac{m\hbar\omega}{2} (a_2^\dagger - a_2)^2 = -\frac{m\hbar\omega}{2} ((a_2^\dagger)^2 + a_2^2 - a_2^\dagger a_2 - a_2 a_2^\dagger). \]
    Since $a_2 a_2^\dagger = N_2 + 1$, the term $-(a_2^\dagger a_2 + a_2 a_2^\dagger) = -(N_2 + N_2 + 1) = -(2N_2 + 1)$. Thus:
    \[ p_2^2 = \frac{m\hbar\omega}{2} (2N_2 + 1 - a_2^2 - (a_2^\dagger)^2). \]
\end{itemize}
The full perturbation operator $V = \lambda x_1^2 p_2^2$ is:
\[ V = \lambda \left( \frac{\hbar}{2m\omega} \right) \left( \frac{m\hbar\omega}{2} \right) (a_1^2 + (a_1^\dagger)^2 + 2N_1 + 1) (2N_2 + 1 - a_2^2 - (a_2^\dagger)^2). \]
\[ V = \frac{\lambda\hbar^2}{4} \underbrace{(a_1^2 + (a_1^\dagger)^2 + 2N_1 + 1)}_{A} \underbrace{(2N_2 + 1 - a_2^2 - (a_2^\dagger)^2)}_{B}. \]

\textbf{3. Calculating Matrix Elements $V_{ij} = \langle \psi_i | V | \psi_j \rangle$:}

We calculate the matrix in the basis $\{|2,0\rangle, |1,1\rangle, |0,2\rangle\}$. 
\sidenote{In degenerate perturbation theory, the first-order energy shifts are the eigenvalues of the perturbation matrix $V$ within the degenerate subspace. For any two states $|\psi_i\rangle, |\psi_j\rangle$ in this subspace (where $n_1+n_2 = 2$), the matrix element $\langle \psi_i | V | \psi_j \rangle$ is non-zero \textbf{only if} $V$ contains terms that "connect" these states while keeping the total number of quanta constant. Any term in $V$ that changes the total $n$ (like $a_1^2$, which gives $n \to n-2$) will produce a state orthogonal to the subspace, resulting in a zero inner product.}

\begin{itemize}
    \item \textbf{Calculating $V_{11} = \langle 2, 0 | V | 2, 0 \rangle$:}
    For diagonal elements, we only need terms in the product $AB$ that leave the state $|n_1, n_2\rangle$ unchanged. These are the \textbf{diagonal-preserving terms} $(2N_1 + 1)(2N_2 + 1)$. Any other term (like $a_1^2 a_2^2$) would change the state and yield a zero inner product due to orthogonality.
    \[ V_{11} = \frac{\lambda\hbar^2}{4} \langle 2,0 | (2N_1 + 1)(2N_2 + 1) | 2,0 \rangle = \frac{\lambda\hbar^2}{4} (2(2)+1)(2(0)+1) = \frac{5}{4}\lambda\hbar^2. \]

    \item \textbf{Calculating $V_{22} = \langle 1, 1 | V | 1, 1 \rangle$:}
    \[ V_{22} = \frac{\lambda\hbar^2}{4} \langle 1,1 | (2N_1 + 1)(2N_2 + 1) | 1,1 \rangle = \frac{\lambda\hbar^2}{4} (2(1)+1)(2(1)+1) = \frac{9}{4}\lambda\hbar^2. \]

    \item \textbf{Calculating $V_{33} = \langle 0, 2 | V | 0, 2 \rangle$:}
    \[ V_{33} = \frac{\lambda\hbar^2}{4} \langle 0,2 | (2N_1 + 1)(2N_2 + 1) | 0,2 \rangle = \frac{\lambda\hbar^2}{4} (2(0)+1)(2(2)+1) = \frac{5}{4}\lambda\hbar^2. \]

    \item \textbf{Calculating $V_{13} = \langle 2, 0 | V | 0, 2 \rangle$:}
    We need terms that can convert $|0, 2\rangle$ into $|2, 0\rangle$. The only term is $(a_1^\dagger)^2 (-a_2^2)$.
    \[ V_{13} = \frac{\lambda\hbar^2}{4} \langle 2,0 | -(a_1^\dagger)^2 a_2^2 | 0,2 \rangle. \]
    Action: $a_2^2 |0, 2\rangle = \sqrt{2 \cdot 1} |0, 0\rangle = \sqrt{2} |0, 0\rangle$. 
    Then $-(a_1^\dagger)^2 (\sqrt{2} |0, 0\rangle) = -\sqrt{2} (\sqrt{1 \cdot 2} |2, 0\rangle) = -2 |2, 0\rangle$.
    \[ V_{13} = \frac{\lambda\hbar^2}{4} (-2) = -\frac{1}{2}\lambda\hbar^2. \]
    By Hermiticity, $V_{31} = V_{13} = -\frac{1}{2}\lambda\hbar^2$.

    \item \textbf{Calculating $V_{12}$ and $V_{23}$:}
    These vanish because any term in the product $AB$ (like $a_1^2 (2N_2+1)$) changes $n_1$ by 2 but does not compensate with $n_2$. There are no terms that change $(n_1, n_2)$ from $(2,0) \to (1,1)$.
\end{itemize}

\textbf{4. Diagonalizing the Matrix:}
The perturbation matrix $W$ is:
\[ W = \frac{\lambda\hbar^2}{4} \begin{pmatrix} 5 & 0 & -2 \\ 0 & 9 & 0 \\ -2 & 0 & 5 \end{pmatrix}. \]
The eigenvalue for $|1, 1\rangle$ is already visible as \textbf{$E^{(1)}_2 = \frac{9}{4}\lambda\hbar^2$}.
For the remaining $2 \times 2$ block:
\[ \det \left[ \frac{\lambda\hbar^2}{4} \begin{pmatrix} 5 & -2 \\ -2 & 5 \end{pmatrix} - \epsilon I \right] = 0. \]
Let $x = \epsilon / (\lambda\hbar^2 / 4)$. Then $(5-x)^2 - (-2)^2 = 0$:
\[ (5-x)^2 = 4 \implies 5-x = \pm 2. \]
\begin{itemize}
    \item $5-x = 2 \implies x = 3 \implies \mathbf{E^{(1)}_1 = \frac{3}{4}\lambda\hbar^2}$.
    \item $5-x = -2 \implies x = 7 \implies \mathbf{E^{(1)}_3 = \frac{7}{4}\lambda\hbar^2}$.
\end{itemize}

\textbf{Conclusion:}
The first-order energy shifts for the $n=2$ level are $\left\{ \frac{3}{4}\lambda\hbar^2, \frac{7}{4}\lambda\hbar^2, \frac{9}{4}\lambda\hbar^2 \right\}$.

\begin{remark}
\textbf{Full Expansion of the Perturbation $V$:}
To see why $(2N_1+1)(2N_2+1)$ is the only diagonal term, we expand the product $V = \frac{\lambda\hbar^2}{4} AB$:
\[ \begin{aligned}
V = \frac{\lambda\hbar^2}{4} \Big[ & \underbrace{(2N_1+1)(2N_2+1)}_{\text{Diagonal: preserves both } n_1, n_2} \\
& \underbrace{- (2N_1+1)a_2^2 - (2N_1+1)(a_2^\dagger)^2}_{\text{Changes } n_2 \text{ by } \pm 2, \text{ changes total quanta } n} \\
& \underbrace{+ a_1^2(2N_2+1) + (a_1^\dagger)^2(2N_2+1)}_{\text{Changes } n_1 \text{ by } \pm 2, \text{ changes total quanta } n} \\
& \underbrace{- a_1^2 a_2^2 - (a_1^\dagger)^2 (a_2^\dagger)^2}_{\text{Changes total quanta } n \text{ by } \pm 4} \\
& \underbrace{- a_1^2 (a_2^\dagger)^2 - (a_1^\dagger)^2 a_2^2}_{\text{Preserves total } n, \text{ shifts } 2 \text{ quanta between oscillators}} \Big]
\end{aligned} \]
When calculating the diagonal elements $\langle n_1, n_2 | V | n_1, n_2 \rangle$, any term that changes $n_1$ or $n_2$ will result in a state orthogonal to the original basis state, yielding zero. Thus, only the first term contributes. For the off-diagonal element $V_{13} = \langle 2, 0 | V | 0, 2 \rangle$, the only term that can "convert" two quanta from oscillator 2 into oscillator 1 is the very last term $-(a_1^\dagger)^2 a_2^2$.
\end{remark}
\end{solution}

\begin{example}[Common Lie Algebras in Physics]
In mathematical physics, several Lie algebras appear frequently due to their relationship with symmetries. The commutation relations for the basis elements define the algebra's structure:

\begin{enumerate}
    \item \textbf{Angular Momentum and $su(2)$:}
    The components of angular momentum $\mathbf{L}$ satisfy:
    \[ [L_i, L_j] = \mathrm{i} \hbar \sum_{k=1}^3 \epsilon_{ijk} L_k \]
    Using the \textbf{Pauli matrices} $\sigma_i$, a representation is $J_i = \frac{\hbar}{2}\sigma_i$. The ladder operators $J_\pm = J_x \pm \mathrm{i}J_y$ (as used in Exercise 2) satisfy $[J_z, J_\pm] = \pm \hbar J_\pm$.
    
    \item \textbf{The Heisenberg Algebra $h(1)$:}
    Defined by the canonical commutation relations (CCR) between position and momentum:
    \[ [x, p] = \mathrm{i} \hbar I, \quad [x, I] = 0, \quad [p, I] = 0 \]
    where $I$ is the identity operator.
    
    \item \textbf{The Lorentz Algebra $so(3,1)$:}
    Describes spacetime rotations and boosts. The generators of rotations $J_i$ and boosts $K_i$ satisfy:
    \[ [J_i, J_j] = \mathrm{i} \epsilon_{ijk} J_k, \quad [J_i, K_j] = \mathrm{i} \epsilon_{ijk} K_k, \quad [K_i, K_j] = -\mathrm{i} \epsilon_{ijk} J_k \]
\end{enumerate}
\end{example}

\begin{exercise}
\label{physics:2022:3}
A Killing vector field $k^\mu \frac{\partial}{\partial x^\mu}$ satisfies the equation $k^\lambda \partial_\lambda g_{\mu\nu} + \partial_\mu k^\lambda g_{\lambda\nu} + \partial_\nu k^\lambda g_{\lambda\mu} = 0$.\sidenote{A \textbf{Killing vector field} is a vector field on a Riemannian or pseudo-Riemannian manifold that preserves the metric. It represents a \textbf{generator of an isometry} (a symmetry of the spacetime). In physics, via \textbf{Noether's Theorem}, each Killing vector corresponds to a conserved quantity, such as energy (if there is a time-like Killing vector) or momentum (if there is a space-like Killing vector).}

\begin{enumerate}[label=(\alph*)]
    \item Prove: $\nabla_\mu k_\nu + \nabla_\nu k_\mu = 0$, here $\nabla_\mu$ is the covariant derivative.
    \item For a moving particle in gravitational background with a Killing vector field, prove: $k^\mu P_\mu$ is a conserved quantity. Here $P_\mu = m \frac{dx^\nu}{d\tau}g_{\mu\nu}$ is the momentum for the free falling particle with trajectory $x^\nu(\tau)$.
\end{enumerate}
\end{exercise}

\begin{solution}
\textbf{Noether's Theorem:}
\begin{theorem}
Every differentiable symmetry of the action of a physical system with conservative forces has a corresponding conservation law. Specifically, if the action is invariant under a continuous transformation generated by a vector field $k^\mu$, there exists a conserved quantity $Q$ along the trajectory of the system ($dQ/d\tau = 0$).
\end{theorem}

\begin{definition}
The \textbf{Killing vector space} of a manifold $(M, g)$ is the set of all vector fields $X$ whose flow consists of isometries.\sidenote{The \textbf{flow} of a vector field $X$ is a family of maps $\phi_t: M \to M$ that "slide" the manifold along the vector field. Specifically, for any point $p$, the curve $\gamma(t) = \phi_t(p)$ is an \textbf{integral curve} whose velocity vector is exactly $X(\gamma(t))$. If the flow consists of \textbf{isometries}, it means that for any $t$, the map $\phi_t$ preserves the metric: $\phi_t^* g = g$. This is the geometric way of saying the spacetime "looks the same" when shifted along $X$.} Mathematically, $X$ is a \textbf{Killing vector field} if the Lie derivative of the metric tensor $g$ with respect to $X$ vanishes:
\[ \mathcal{L}_X g = 0. \]
In local coordinates, this condition is expressed by the \textbf{Killing equation}:
\[ \nabla_\mu X_\nu + \nabla_\nu X_\mu = 0. \]
The set of all Killing vector fields forms a finite-dimensional \textbf{Lie algebra} under the standard Lie bracket of vector fields.
\end{definition}

\textbf{Prerequisites / Core Concepts:}
\begin{itemize}
    \item \textbf{Metric Tensor vs. Riemannian Metric:}
    While often used interchangeably, there is a technical distinction:
    \begin{itemize}
        \item \textbf{Metric Tensor}: The general mathematical object—a symmetric, non-degenerate $(0,2)$-tensor field $g_{\mu\nu}$ that defines the inner product on each tangent space.
        \item \textbf{Riemannian Metric}: A metric tensor that is \textbf{positive-definite} (all eigenvalues are positive). In this geometry, the distance $ds^2$ is always $\geq 0$.
        \item \textbf{Lorentzian (Pseudo-Riemannian) Metric}: The type used in General Relativity (and this exercise). It has a signature like $(-,+,+,+)$. This allows $ds^2$ to be negative (time-like), zero (light-like), or positive (space-like), which is essential for describing the \textbf{causal structure} of spacetime.
    \end{itemize}

    \item \textbf{Deriving the Metric Matrix ($g_{\mu\nu}$):}
    The metric tensor represents the inner product of the coordinate basis vectors $\mathbf{e}_\mu = \frac{\partial \mathbf{R}}{\partial x^\mu}$. If we know the transformation from Cartesian coordinates $(X, Y, Z)$ to our general coordinates $(x^1, x^2, x^3)$, the metric is computed as:
    \[ g_{\mu\nu} = \frac{\partial X}{\partial x^\mu} \frac{\partial X}{\partial x^\nu} + \frac{\partial Y}{\partial x^\mu} \frac{\partial Y}{\partial x^\nu} + \frac{\partial Z}{\partial x^\mu} \frac{\partial Z}{\partial x^\nu}. \]
    \begin{itemize}
        \item \textbf{Example (Spherical)}: With $X = r \sin\theta \cos\phi, Y = r \sin\theta \sin\phi, Z = r \cos\theta$:
        \begin{enumerate}
            \item For $g_{rr}$: $(\partial_r X)^2 + (\partial_r Y)^2 + (\partial_r Z)^2 = (\sin\theta\cos\phi)^2 + (\sin\theta\sin\phi)^2 + (\cos\theta)^2 = 1$.
            \item For $g_{\theta\theta}$: $(\partial_\theta X)^2 + (\partial_\theta Y)^2 + (\partial_\theta Z)^2 = (r\cos\theta\cos\phi)^2 + (r\cos\theta\sin\phi)^2 + (-r\sin\theta)^2 = r^2$.
        \end{enumerate}
        This is how the "scale factors" in the metric are rigorously derived from the geometry of the transformation.
    \end{itemize}

    \item \textbf{The Metric and its Inverse ($g_{\mu\nu}$ and $g^{\mu\nu}$):}
    The metric tensor $g_{\mu\nu}$ (with lower indices) is used to calculate distances and lower indices. Its \textbf{inverse matrix} is denoted $g^{\mu\nu}$ (with upper indices) and is used to raise indices. They are related by the \textbf{orthogonality condition}:
    \[ g^{\mu\lambda} g_{\lambda\nu} = \delta^\mu_\nu, \]
    where $\delta^\mu_\nu$ is the Kronecker delta ($1$ if $\mu=\nu$ and $0$ otherwise). 
    \begin{itemize}
        \item \textbf{Example}: In our 3D spherical metric, $g_{rr} = 1$ and $g_{\theta\theta} = r^2$. Since the matrix is diagonal, the inverse components are simply the reciprocals:
        \[ g^{rr} = \frac{1}{g_{rr}} = 1, \quad g^{\theta\theta} = \frac{1}{g_{\theta\theta}} = \frac{1}{r^2}. \]
        \item \textbf{Off-diagonal terms}: If the coordinates are orthogonal (like $r$ and $\theta$), then $g_{r\theta} = g_{\theta r} = 0$, which implies $g^{r\theta} = 0$ as well.
    \end{itemize}

    \item \textbf{Einstein Summation Convention:} A notation convention where an index that appears twice in a single term (once up and once down) is automatically summed over all its possible values. For example, in 4D spacetime:
    \[ A^\mu B_\mu = \sum_{\mu=0}^3 A^\mu B_\mu = A^0 B_0 + A^1 B_1 + A^2 B_2 + A^3 B_3. \]
    This convention is used throughout this proof for indices like $\mu, \nu, \lambda$.

    \item \textbf{The Lie Derivative ($\mathcal{L}_X$):}
    The Lie derivative of a tensor field $T$ with respect to a vector field $X$ measures the rate of change of $T$ along the \textbf{flow} $\phi_t$ of $X$. It is defined as:
    \[ \mathcal{L}_X T = \lim_{t \to 0} \frac{\phi_t^* T - T}{t}. \]
    For the metric tensor $g_{\mu\nu}$, the Lie derivative in local coordinates is:
    \[ \mathcal{L}_X g_{\mu\nu} = X^\lambda \partial_\lambda g_{\mu\nu} + g_{\lambda\nu} \partial_\mu X^\lambda + g_{\mu\lambda} \partial_\nu X^\lambda. \]
    Using the covariant derivative $\nabla$, this simplifies beautifully to the Killing operator:
    \[ \mathcal{L}_X g_{\mu\nu} = \nabla_\mu X_\nu + \nabla_\nu X_\mu. \]

    \item \textbf{Covariant Derivative and the Connection ($\nabla$):}
    The \textbf{connection} $\nabla$ is the overall operator. We use different notations depending on the direction of differentiation:
    \begin{itemize}
        \item \textbf{$\nabla_X$}: This is the derivative in the direction of a general \textbf{vector field} $X$.
        \item \textbf{$\nabla_\mu$}: This is a shorthand for $\nabla_{\mathbf{e}_\mu}$, the derivative in the direction of the $\mu$-th \textbf{basis vector}.
        \item \textbf{$\nabla_1, \nabla_2, \dots$}: These are specific coordinate derivatives (e.g., $\nabla_1 = \nabla_r$ in spherical coordinates).
    \end{itemize}
    They are related by \textbf{linearity}: If $X = X^\mu \mathbf{e}_\mu$, then:
    \[ \nabla_X = X^\mu \nabla_\mu = X^1 \nabla_1 + X^2 \nabla_2 + \dots \]
    This means the derivative along any path is just a weighted sum of the derivatives along the coordinate axes. It satisfies the \textbf{Leibniz rule}:
    \[ \nabla_X (f Y) = X(f) Y + f \nabla_X Y, \]
    where $f$ is a scalar function. In local coordinates, the components $\nabla_\mu$\sidenote{The \textbf{partial derivative} $\partial_\mu$ only measures how the numerical components of a tensor change. The \textbf{covariant derivative} $\nabla_\mu$ measures how the tensor \textbf{itself} changes, which includes the fact that the \textbf{basis vectors} ($\mathbf{e}_\nu$) are also rotating and stretching as you move through curved space. Thus, $\nabla_\mu = \partial_\mu + (\text{correction for basis change})$.} act on a covector $X_\nu$ as:
    \[ \nabla_\mu X_\nu = \partial_\mu X_\nu - \Gamma^\lambda_{\mu\nu} X_\lambda. \]
    The \textbf{Christoffel symbols} $\Gamma^\lambda_{\mu\nu}$ are fundamentally defined by the action of the connection on the basis vectors:
    \[ \nabla_\mu \mathbf{e}_\nu = \Gamma^\lambda_{\mu\nu} \mathbf{e}_\lambda. \]
    This equation shows that $\Gamma$ tells us how the basis vector $\mathbf{e}_\nu$ "tilts" or "grows" as we move in the $\mu$-direction. For the \textbf{Levi-Civita connection}, they are determined entirely by the metric $g$:
    \[ \Gamma^\rho_{\mu\nu} = \frac{1}{2} g^{\rho\sigma} \left( \partial_\mu g_{\nu\sigma} + \partial_\nu g_{\mu\sigma} - \partial_\sigma g_{\mu\nu} \right). \]
    Geometrically, $\nabla$ defines \textbf{parallel transport}:\sidenote{A vector field $V$ is said to be \textbf{parallel} if its covariant derivative vanishes everywhere: $\nabla V = 0$ (or $\nabla_\mu V^\nu = 0$). 
    If we have two vectors $v \in T_p M$ and $w \in T_q M$ at different points, they are defined as \textbf{parallel} relative to a curve $\gamma$ connecting $p$ and $q$ if $w$ is the \textbf{parallel transport} of $v$ along that curve. In curved spacetime, "parallel" is generally path-dependent.} it allows us to compare tensors at different points by "sliding" them along curves while keeping them "constant" relative to the manifold's curvature. A tensor $T$ is parallelly transported along a curve with velocity $u^\mu$ if $u^\nu \nabla_\nu T = 0$.

    \item \textbf{Geometric Interpretation of the Minus Sign}:
    Why does $\nabla_\mu X_\nu$ have a \textbf{minus} sign while $\nabla_\mu X^\nu$ has a \textbf{plus} sign? 
    A covector $X_\nu$ is a linear map that eats a vector $Y^\nu$ to produce a scalar $S = X_\nu Y^\nu$. Because a scalar is just a number, its derivative is simply the partial derivative: $\nabla_\mu S = \partial_\mu S$.
    Using the \textbf{Leibniz rule}:
    \[ \partial_\mu (X_\nu Y^\nu) = (\nabla_\mu X_\nu) Y^\nu + X_\nu (\nabla_\mu Y^\nu). \]
    We know that for a vector, $\nabla_\mu Y^\nu = \partial_\mu Y^\nu + \Gamma^\nu_{\mu\lambda} Y^\lambda$. Substituting this in:
    \[ \partial_\mu X_\nu Y^\nu + X_\nu \partial_\mu Y^\nu = (\nabla_\mu X_\nu) Y^\nu + X_\nu (\partial_\mu Y^\nu + \Gamma^\nu_{\mu\lambda} Y^\lambda). \]
    Canceling the $X_\nu \partial_\mu Y^\nu$ terms from both sides and rearranging for $\nabla_\mu X_\nu$:
    \[ \nabla_\mu X_\nu = \partial_\mu X_\nu - \Gamma^\lambda_{\mu\nu} X_\lambda. \]
    Geometrically, this means that if the coordinate grid "bends" in one direction (positive $\Gamma$), the \textbf{dual grid} (the covectors) must "bend" in the \textbf{opposite direction} to keep the measurements (scalar products) consistent. The minus sign is the mathematical "correction" that maintains the duality.
    \item \textbf{Metric Compatibility ($\nabla_\lambda g_{\mu\nu} = 0$):}
    The defining property of the Levi-Civita connection is that it "preserves" the metric. To see why this implies the formula for $\partial_\lambda g_{\mu\nu}$, recall the general rule for the covariant derivative of a \textbf{rank-2 tensor} $T_{\mu\nu}$:\sidenote{A \textbf{$(p,q)$ tensor} $T^{\alpha_1 \dots \alpha_p}_{\beta_1 \dots \beta_q}$ is a geometric object with $p$ upper indices (contravariant) and $q$ lower indices (covariant). The general rule for its \textbf{covariant derivative} is:
    \[ \begin{aligned}
    \nabla_\gamma T^{\alpha_1 \dots \alpha_p}_{\beta_1 \dots \beta_q} &= \partial_\gamma T^{\alpha_1 \dots \alpha_p}_{\beta_1 \dots \beta_q} \\
    &+ \sum_{i=1}^p \Gamma^{\alpha_i}_{\gamma \sigma} T^{\dots \sigma \dots}_{\dots} \quad (\text{positive } \Gamma \text{ for each upper index}) \\
    &- \sum_{j=1}^q \Gamma^{\sigma}_{\gamma \beta_j} T^{\dots}_{\dots \sigma \dots} \quad (\text{negative } \Gamma \text{ for each lower index}).
    \end{aligned} \]
    For the metric $g_{\mu\nu}$, we have $(p,q)=(0,2)$, so we get exactly two negative Christoffel terms and zero positive ones.}
    \[ \nabla_\lambda T_{\mu\nu} = \partial_\lambda T_{\mu\nu} - \Gamma^\rho_{\lambda\mu} T_{\rho\nu} - \Gamma^\rho_{\lambda\nu} T_{\mu\rho}. \]
    By substituting the metric tensor $g_{\mu\nu}$ for $T_{\mu\nu}$, we have:
    \[ \nabla_\lambda g_{\mu\nu} = \partial_\lambda g_{\mu\nu} - \Gamma^\rho_{\lambda\mu} g_{\rho\nu} - \Gamma^\rho_{\lambda\nu} g_{\mu\rho}. \]
    Setting $\nabla_\lambda g_{\mu\nu} = 0$ (the compatibility condition) and moving the Christoffel terms to the other side of the equation gives:
    \[ \partial_\lambda g_{\mu\nu} = \Gamma^\rho_{\lambda\mu} g_{\rho\nu} + \Gamma^\rho_{\lambda\nu} g_{\mu\rho}. \]
    This identity is the crucial "bridge" used in part (a) to convert coordinate derivatives into geometric ones. It ensures that lengths and angles are preserved under parallel transport.
    \item \textbf{Geodesic Equation:} The equation of motion for a free-falling particle in curved spacetime: $u^\nu \nabla_\nu u^\mu = 0$, where $u^\mu = \frac{dx^\mu}{d\tau}$ is the four-velocity.
\end{itemize}

\begin{remark}[3D Example: Spherical Coordinates]
To understand $\nabla_\lambda$, let's look at 3D space ($D=3$) in \textbf{spherical coordinates} $(x^1, x^2, x^3) = (r, \theta, \phi)$.
\begin{enumerate}
    \item \textbf{The Index $\lambda$}: The index $\lambda$ is a label for which coordinate direction we are moving in. Here, $\lambda \in \{r, \theta, \phi\}$.
    \item \textbf{The Operator $\nabla_\lambda$}: $\nabla_r$ is the derivative along the radius, $\nabla_\theta$ along the latitude, etc.
    \item \textbf{The Metric $g_{\mu\nu}$}: In these coordinates, the distance $ds^2 = dr^2 + r^2 d\theta^2 + r^2 \sin^2\theta d\phi^2$ gives:
    \[ g_{\mu\nu} = \begin{pmatrix} 1 & 0 & 0 \\ 0 & r^2 & 0 \\ 0 & 0 & r^2 \sin^2\theta \end{pmatrix}. \]
    \item \textbf{Christoffel Symbols $\Gamma$}: Since the basis vectors change direction as we move (e.g., $\hat{\theta}$ depends on position), the partial derivative is not enough. For example, the symbol $\Gamma^r_{\theta\theta}$ is calculated as:
    \[ \Gamma^r_{\theta\theta} = \frac{1}{2} g^{rr} (\partial_\theta g_{\theta r} + \partial_\theta g_{r \theta} - \partial_r g_{\theta\theta}) = \frac{1}{2} (1) (0 + 0 - 2r) = -r. \]
    \item \textbf{The Covariant Derivative $\nabla_\theta X^r$}: Suppose we have a vector field $X$. The derivative of its radial component in the $\theta$ direction is:
    \[ \nabla_\theta X^r = \partial_\theta X^r + \Gamma^r_{\theta\lambda} X^\lambda = \partial_\theta X^r + \Gamma^r_{\theta\theta} X^\theta = \partial_\theta X^r - r X^\theta. \]
    The term $-r X^\theta$ is the "correction" that accounts for the fact that the radial direction itself rotates as you move along $\theta$.
    \item \textbf{Geodesic Equation}: The radial ($r$) component of the geodesic equation is:
    \[ \frac{d^2 r}{d\tau^2} + \Gamma^r_{\theta\theta} \left(\frac{d\theta}{d\tau}\right)^2 + \Gamma^r_{\phi\phi} \left(\frac{d\phi}{d\tau}\right)^2 = 0 \implies \frac{d^2 r}{d\tau^2} - r \dot{\theta}^2 - r \sin^2\theta \dot{\phi}^2 = 0. \]
    This is exactly $a_r = 0$ in polar coordinates, where the $\Gamma$ terms represent the \textbf{centrifugal acceleration}.
\end{enumerate}
\end{remark}

\begin{example}[Consistency of Inner Products]
Metric compatibility $\nabla_\lambda g_{\mu\nu} = 0$ ensures that the \textbf{inner product} of two parallel-transported vectors is constant along a curve. Consider two vectors $V$ and $W$ parallelly transported along a curve with velocity $u^\lambda$ (so $u^\lambda \nabla_\lambda V^\mu = 0$ and $u^\lambda \nabla_\lambda W^\nu = 0$). The rate of change of their inner product is:
\[ \nabla_u (g_{\mu\nu} V^\mu W^\nu) = (\nabla_u g_{\mu\nu}) V^\mu W^\nu + g_{\mu\nu} (\nabla_u V^\mu) W^\nu + g_{\mu\nu} V^\mu (\nabla_u W^\nu). \]
The last two terms vanish by the parallel transport condition. The first term is:
\[ (\nabla_u g_{\mu\nu}) V^\mu W^\nu = (u^\lambda \nabla_\lambda g_{\mu\nu}) V^\mu W^\nu = 0 \quad (\text{by metric compatibility}). \]
Thus, $g_{\mu\nu} V^\mu W^\nu$ is constant. This means the \textbf{lengths} of vectors and the \textbf{angles} between them are preserved during parallel transport.
\end{example}

\begin{example}[Geodesics in Flat Space]
In flat Euclidean space using Cartesian coordinates, the metric is constant ($g_{\mu\nu} = \eta_{\mu\nu}$) and all Christoffel symbols vanish: $\Gamma^\rho_{\mu\nu} = 0$. The geodesic equation:
\[ \frac{d u^\mu}{d\tau} + \Gamma^\mu_{\nu\rho} u^\nu u^\rho = 0 \]
reduces to:
\[ \frac{d^2 x^\mu}{d\tau^2} = 0. \]
The solution is a \textbf{straight line} $x^\mu(\tau) = v^\mu \tau + x_0^\mu$. Thus, the geodesic equation is the formal generalization of "straight line motion" to curved manifolds, where the particle follows the "straightest possible path" allowed by the geometry.
\end{example}

\begin{example}[Full Computation of Christoffel Symbols]
In 3D spherical coordinates $(r, \theta, \phi)$, the non-zero metric components are $g_{rr}=1, g_{\theta\theta}=r^2, g_{\phi\phi}=r^2\sin^2\theta$. The inverse components are $g^{rr}=1, g^{\theta\theta}=1/r^2, g^{\phi\phi}=1/(r^2\sin^2\theta)$.
We compute the non-zero derivatives:
\[ \partial_r g_{\theta\theta} = 2r, \quad \partial_r g_{\phi\phi} = 2r \sin^2 \theta, \quad \partial_\theta g_{\phi\phi} = 2r^2 \sin\theta\cos\theta. \]
Using the formula $\Gamma^\rho_{\mu\nu} = \frac{1}{2} g^{\rho\sigma} (\partial_\mu g_{\nu\sigma} + \partial_\nu g_{\mu\sigma} - \partial_\sigma g_{\mu\nu})$:
\begin{enumerate}
    \item \textbf{Calculating $\Gamma^r_{\theta\theta}$}:
    \[ \Gamma^r_{\theta\theta} = \frac{1}{2} g^{rr} (\partial_\theta g_{\theta r} + \partial_\theta g_{r\theta} - \partial_r g_{\theta\theta}) = \frac{1}{2} (1) (0 + 0 - 2r) = -r. \]
    
    \item \textbf{Calculating $\Gamma^\theta_{r\theta}$}:
    \[ \Gamma^\theta_{r\theta} = \frac{1}{2} g^{\theta\theta} (\partial_r g_{\theta\theta} + \partial_\theta g_{r\theta} - \partial_\theta g_{r\theta}) = \frac{1}{2} \left(\frac{1}{r^2}\right) (2r + 0 - 0) = \frac{1}{r}. \]
    
    \item \textbf{Calculating $\Gamma^\theta_{\phi\phi}$}:
    \[ \Gamma^\theta_{\phi\phi} = \frac{1}{2} g^{\theta\theta} (\partial_\phi g_{\phi\theta} + \partial_\phi g_{\theta\phi} - \partial_\theta g_{\phi\phi}) = \frac{1}{2} \left(\frac{1}{r^2}\right) (0 + 0 - 2r^2 \sin\theta\cos\theta) = -\sin\theta\cos\theta. \]
\end{enumerate}
\end{example}

\textbf{Intuition / Motivation:}
A Killing vector points in a "direction of symmetry" for the spacetime. For example, if a spacetime is perfectly spherical, rotating around it shouldn't change the physics. The Killing equation mathematically encodes this "no change" condition. Because Noether's theorem tells us that symmetries lead to conservation laws, moving along a geodesic in a spacetime with a symmetry implies the momentum in the direction of that symmetry is strictly conserved. We will prove this by showing the derivative of the momentum along the path is exactly zero.

\textbf{Logical Step-by-Step Breakdown:}

\textbf{(a) Proof of the Covariant Killing Equation:}
\begin{enumerate}
    \item \textbf{Step 1: Using Metric Compatibility as a Bridge.}
    The original Killing equation is given in terms of partial derivatives $\partial_\lambda g_{\mu\nu}$. To move to a covariant form, we must replace these partial derivatives with Christoffel symbols. This is where \textbf{metric compatibility} ($\nabla_\lambda g_{\mu\nu} = 0$) is essential. 
    By the definition of the covariant derivative of a $(0,2)$ tensor:
    \[ \nabla_\lambda g_{\mu\nu} = \partial_\lambda g_{\mu\nu} - \Gamma^\rho_{\lambda\mu} g_{\rho\nu} - \Gamma^\rho_{\lambda\nu} g_{\mu\rho}. \]
    Since the Levi-Civita connection is metric-compatible ($\nabla_\lambda g_{\mu\nu} = 0$), we can solve for the partial derivative:
    \[ \partial_\lambda g_{\mu\nu} = \Gamma^\rho_{\lambda\mu} g_{\rho\nu} + \Gamma^\rho_{\lambda\nu} g_{\mu\rho}. \]
    This substitution allows us to translate the "flat" partial derivatives of the metric into the "curved" geometric connection.

    \begin{definition}
    A \textbf{$(p,q)$ tensor} $T^{\alpha_1 \dots \alpha_p}_{\beta_1 \dots \beta_q}$ has $p$ contravariant (upper) indices and $q$ covariant (lower) indices. Its \textbf{covariant derivative} $\nabla_\gamma$ is defined by the following expansion:
    \[ \begin{aligned}
    \nabla_\gamma T^{\alpha_1 \dots \alpha_p}_{\beta_1 \dots \beta_q} &= \partial_\gamma T^{\alpha_1 \dots \alpha_p}_{\beta_1 \dots \beta_q} \\
    &+ \sum_{i=1}^p \Gamma^{\alpha_i}_{\gamma \sigma} T^{\alpha_1 \dots \sigma \dots \alpha_p}_{\beta_1 \dots \beta_q} \quad (\text{one } +\Gamma \text{ term for each upper index}) \\
    &- \sum_{j=1}^q \Gamma^{\sigma}_{\gamma \beta_j} T^{\alpha_1 \dots \alpha_p}_{\beta_1 \dots \sigma \dots \beta_q} \quad (\text{one } -\Gamma \text{ term for each lower index}).
    \end{aligned} \]
    For a $(0,2)$ tensor like the metric $g_{\mu\nu}$, this reduces to the two negative terms shown in Step 1.
    \end{definition}

    \begin{remark}[Commutativity of Metric and Connection]
    The identity $g_{\lambda\nu} \nabla_\mu k^\lambda = \nabla_\mu k_\nu$ is a direct consequence of \textbf{metric compatibility}. To prove it, we start with the definition of the covector $k_\nu = g_{\lambda\nu} k^\lambda$ and differentiate both sides:
    \[ \nabla_\mu k_\nu = \nabla_\mu (g_{\lambda\nu} k^\lambda). \]
    Applying the \textbf{Leibniz rule} (the product rule) for covariant derivatives:
    \[ \nabla_\mu (g_{\lambda\nu} k^\lambda) = \underbrace{(\nabla_\mu g_{\lambda\nu})}_{=0} k^\lambda + g_{\lambda\nu} (\nabla_\mu k^\lambda). \]
    Since the Levi-Civita connection is metric-compatible ($\nabla_\mu g_{\lambda\nu} = 0$), the first term vanishes identically. Thus:
    \[ \nabla_\mu k_\nu = g_{\lambda\nu} \nabla_\mu k^\lambda. \]
    This means the metric can "pass through" the covariant derivative operator. In physics terms, \textbf{lowering an index and then differentiating is the same as differentiating and then lowering the index}.
    \end{remark}
    
    \item \textbf{Step 2: Express partial derivatives of the vector components.}
    Similarly, for the vector field $k^\lambda$:
    \[ \nabla_\mu k^\lambda = \partial_\mu k^\lambda + \Gamma^\lambda_{\mu\sigma} k^\sigma \implies \partial_\mu k^\lambda = \nabla_\mu k^\lambda - \Gamma^\lambda_{\mu\sigma} k^\sigma. \]
    
    \item \textbf{Step 3: Substitute and Simplify.}
    Substitute these expressions into the original Killing equation $k^\lambda \partial_\lambda g_{\mu\nu} + (\partial_\mu k^\lambda) g_{\lambda\nu} + (\partial_\nu k^\lambda) g_{\lambda\mu} = 0$:
    \[ \underbrace{k^\lambda ( \Gamma^\rho_{\lambda\mu} g_{\rho\nu} + \Gamma^\rho_{\lambda\nu} g_{\mu\rho} )}_{\text{From } \partial g} + \underbrace{(\nabla_\mu k^\lambda - \Gamma^\lambda_{\mu\sigma} k^\sigma) g_{\lambda\nu}}_{\text{From } \partial k} + \underbrace{(\nabla_\nu k^\lambda - \Gamma^\lambda_{\nu\sigma} k^\sigma) g_{\lambda\mu}}_{\text{From } \partial k} = 0. \]
    
    \item \textbf{Step 4: Distribute indices and cancel.}
    Lowering the indices ($g_{\lambda\nu} \nabla_\mu k^\lambda = \nabla_\mu k_\nu$)\sidenote{Index lowering $k_\nu = g_{\lambda\nu} k^\lambda$ is the fundamental way we define the \textbf{dual covector} associated with a vector $k^\lambda$ using the inner product. Crucially, because $\nabla_\mu g_{\lambda\nu} = 0$, the metric is "constant" relative to the connection. This allows the covariant derivative to \textbf{commute} with the metric: $\nabla_\mu (g_{\lambda\nu} k^\lambda) = (\nabla_\mu g_{\lambda\nu}) k^\lambda + g_{\lambda\nu} (\nabla_\mu k^\lambda) = g_{\lambda\nu} \nabla_\mu k^\lambda$. Thus, differentiating a lowered index is the same as lowering the result of a differentiation.}:
    \[ k^\lambda \Gamma^\rho_{\lambda\mu} g_{\rho\nu} + k^\lambda \Gamma^\rho_{\lambda\nu} g_{\mu\rho} + \nabla_\mu k_\nu - \Gamma^\lambda_{\mu\sigma} k^\sigma g_{\lambda\nu} + \nabla_\nu k_\mu - \Gamma^\lambda_{\nu\sigma} k^\sigma g_{\lambda\mu} = 0. \]
    Because $\Gamma^\rho_{\lambda\mu} = \Gamma^\rho_{\mu\lambda}$ (torsion-free)\sidenote{The \textbf{torsion tensor} is defined as $T(X, Y) = \nabla_X Y - \nabla_Y X - [X, Y]$. For coordinate basis vectors $\mathbf{e}_\mu = \partial_\mu$, the Lie bracket $[\partial_\mu, \partial_\nu]$ is zero. A connection is \textbf{torsion-free} if $T = 0$, which implies:
    \[ \nabla_\mu \mathbf{e}_\nu - \nabla_\nu \mathbf{e}_\mu = 0. \]
    Expanding this using the definition of Christoffel symbols ($\nabla_\mu \mathbf{e}_\nu = \Gamma^\lambda_{\mu\nu} \mathbf{e}_\lambda$):
    \[ (\Gamma^\lambda_{\mu\nu} - \Gamma^\lambda_{\nu\mu}) \mathbf{e}_\lambda = 0 \implies \Gamma^\lambda_{\mu\nu} = \Gamma^\lambda_{\nu\mu}. \]
    This symmetry is what allows us to swap the bottom indices and cancel the terms in this proof.}, the terms $k^\lambda \Gamma^\rho_{\lambda\mu} g_{\rho\nu}$ and $-\Gamma^\lambda_{\mu\sigma} k^\sigma g_{\lambda\nu}$ are identical and cancel out. The same happens for the $\nu$ terms.
    
    \item \textbf{Step 5: Conclusion.}
    We are left with the covariant form:
    \[ \nabla_\mu k_\nu + \nabla_\nu k_\mu = 0. \]
\end{enumerate}

\textbf{(b) Conservation of Momentum along a Geodesic:}
\begin{enumerate}
    \item \textbf{Step 1: Define the conserved quantity.}
    We want to show that the quantity $Q = k^\mu P_\mu$ is conserved along the particle's trajectory $x^\mu(\tau)$.
    The momentum is $P_\mu = m g_{\mu\nu} \frac{dx^\nu}{d\tau} = m u_\mu$, where $u^\mu = \frac{dx^\mu}{d\tau}$ is the four-velocity.
    So, $Q = m k^\mu u_\mu = m k_\mu u^\mu$.
    
    \item \textbf{Step 2: Differentiate $Q$ with respect to proper time $\tau$.}
    The rate of change of $Q$ along the path is given by the directional derivative $u^\nu \nabla_\nu$.
    \[ \frac{dQ}{d\tau} = u^\nu \nabla_\nu (m k_\mu u^\mu) = m u^\nu (\nabla_\nu k_\mu) u^\mu + m k_\mu (u^\nu \nabla_\nu u^\mu) \]
    
    \item \textbf{Step 3: Apply the geodesic equation.}
    For a free-falling particle, the four-velocity satisfies the geodesic equation $u^\nu \nabla_\nu u^\mu = 0$.
    This immediately eliminates the second term:
    \[ \frac{dQ}{d\tau} = m u^\nu u^\mu \nabla_\nu k_\mu \]
    
    \item \textbf{Step 4: Use the Killing equation's antisymmetry.}
    Notice that the tensor $u^\nu u^\mu$ is perfectly symmetric under the exchange of $\mu$ and $\nu$.
    From part (a), the Killing equation $\nabla_\mu k_\nu + \nabla_\nu k_\mu = 0$ means the tensor $\nabla_\nu k_\mu$ is perfectly antisymmetric: $\nabla_\nu k_\mu = -\nabla_\mu k_\nu$.
    When we contract a symmetric tensor with an antisymmetric tensor, the result is identically zero:
    \[ u^\nu u^\mu \nabla_\nu k_\mu = \frac{1}{2} u^\nu u^\mu (\nabla_\nu k_\mu + \nabla_\mu k_\nu) = 0 \]
    
    \item \textbf{Step 5: Conclusion.}
    We have shown that $\frac{dQ}{d\tau} = 0$, which proves that the momentum $k^\mu P_\mu$ projected along the Killing vector field is a conserved quantity for a free-falling particle.
\end{enumerate}
\end{solution}

\begin{exercise}
\label{physics:2022:4}
Consider following metric
\begin{equation}
    ds^2 = -(1 - \frac{2 M}{r}) dv^2 + dv dr + r^2 d\Omega^2
\end{equation}
Here $d\Omega^2$ is the standard metric on two sphere. Consider the hypersurface defined by $S = r - 2M = 0$, and a vector field $l = \tilde{f}(x) (g^{\mu\nu} \partial_\nu S) \frac{\partial}{\partial x^\mu}$\sidenote{In differential geometry, a vector field $l$ is a geometric object expressed in a coordinate basis as $l = l^\mu \frac{\partial}{\partial x^\mu}$ using the \textbf{Einstein summation convention} (where repeated indices like $\mu$ are summed over). Here, $l^\mu = \tilde{f}(x) g^{\mu\nu} \partial_\nu S$ are the \textbf{components} of the vector. We write the full vector by multiplying these coefficients by the basis vectors $\frac{\partial}{\partial x^\mu}$. The gradient form $g^{\mu\nu} \partial_\nu S$ is chosen because the gradient of a function $S$ is always \textbf{orthogonal} to the surfaces where $S$ is constant (like our hypersurface $S=0$).}, here $\tilde{f}(x)$ is a non-zero function. Prove:

\begin{enumerate}[label=(\alph*)]
    \item $l$ is normal to the surface $S$.
    \item $l^2 = 0$ on the surface $S$.
    \item $\frac{\partial}{\partial v}$ is a Killing vector field.
\end{enumerate}
\end{exercise}

\begin{solution}
\textbf{Prerequisites / Core Concepts:}
\begin{itemize}
    \item \textbf{Spacetime Interval ($ds^2$):} Also known as the \textbf{line element}, $ds^2$ defines the differential "distance" (geometry) on a manifold. In this problem, it is defined over a \textbf{4-dimensional} manifold (spacetime) with coordinates $(v, r, \theta, \phi)$, where $v$ is a time-like coordinate and $r, \theta, \phi$ are spatial coordinates.
    \item \textbf{Normal Vector:} A vector field $l^\mu$ is normal to a hypersurface defined by $S(x) = 0$ if it is orthogonal to every vector tangent to the surface. Mathematically, this means $l^\mu$ is proportional to the gradient of the surface, $g^{\mu\nu}\partial_\nu S$.
    \item \textbf{Null Surface:} A surface whose normal vector has a length (norm) of zero, meaning $l_\mu l^\mu = 0$. In general relativity, this characterizes the event horizon of a black hole.
    \item \textbf{Killing Vector Field:} A vector field along which the metric tensor remains invariant. If the metric components are independent of a coordinate $x^\alpha$, then $\frac{\partial}{\partial x^\alpha}$ is a Killing vector.
    \item \textbf{Lie Derivative ($\mathcal{L}_k$):} The fundamental mathematical definition of the Lie derivative of a tensor field $T$ along a vector field $k$ is given by the limit:
    \[ \mathcal{L}_k T = \lim_{t \to 0} \frac{\phi_t^* (T_{\phi_t(p)}) - T_p}{t} \]
    where $\phi_t$ is the \textbf{flow} (one-parameter group of diffeomorphisms) generated by $k$, and $\phi_t^*$ is the \textbf{pull-back} map. Intuitively, it measures the change of the tensor $T$ as it is "dragged" along the integral curves of $k$ and compared to its value at the original point. Unlike the covariant derivative, the Lie derivative does \textbf{not} require a connection or metric; it is a purely differential-topological object. For a $(0,2)$ tensor like the metric $g_{\mu\nu}$, this limit leads to the component formula:
    \[ \mathcal{L}_k g_{\mu\nu} = k^\lambda \partial_\lambda g_{\mu\nu} + g_{\lambda\nu} \partial_\mu k^\lambda + g_{\mu\lambda} \partial_\nu k^\lambda \]
    \item \textbf{Standard Metric on the 2-Sphere ($d\Omega^2$):} The term $d\Omega^2$ is the standard Riemannian metric (line element) on a unit 2-sphere $S^2$. In spherical coordinates $(\theta, \phi)$, with $\theta \in [0, \pi]$ and $\phi \in [0, 2\pi)$, it is given by:
    \[ d\Omega^2 = d\theta^2 + \sin^2\theta d\phi^2. \]
    This defines the intrinsic geometry of the sphere, where $\theta$ is the polar angle and $\phi$ is the azimuthal angle. The associated area element is $dA = \sin\theta d\theta d\phi$.
\end{itemize}

\textbf{Intuition / Motivation:}
We are analyzing the geometry of a black hole using Eddington-Finkelstein coordinates, which are designed to behave smoothly exactly at the event horizon ($r=2M$). Normally, a "surface" like a sphere has a normal vector pointing strictly outward. But at the event horizon of a black hole, space and time are so warped that the "outward" direction becomes a path of light. We show this by proving the normal vector to the horizon has zero length (it is "null"), meaning it points along the trajectory of a photon.

\textbf{Logical Step-by-Step Breakdown:}

\textbf{(a) Normality Proof:}
\begin{enumerate}
    \item \textbf{Step 1: Define a tangent vector.}
    Let $T^\mu$ be any vector tangent to the hypersurface $S(x) = r - 2M = 0$. By definition of a tangent vector, the directional derivative of $S$ along $T^\mu$ must be zero:
    \[ T^\mu \partial_\mu S = 0 \]
    
    \item \textbf{Step 2: Check orthogonality with $l^\mu$.}
    The given vector field is $l^\mu = \tilde{f}(x) g^{\mu\nu} \partial_\nu S$.
    To show $l^\mu$ is normal to the surface, we compute its inner product with the tangent vector $T^\mu$:\sidenote{The inner product (or scalar product) of two vectors $l^\mu$ and $T^\alpha$ in curved spacetime is defined as $g_{\mu\alpha} l^\mu T^\alpha$. The metric tensor $g_{\mu\alpha}$ acts as the "dot product machine," correctly accounting for the geometry. This operation is equivalent to "lowering the index" of one vector (e.g., creating the covariant vector $l_\alpha = g_{\mu\alpha} l^\mu$) and then contracting it with the other ($l_\alpha T^\alpha$).}
    \[ g_{\mu\alpha} l^\mu T^\alpha = g_{\mu\alpha} (\tilde{f}(x) g^{\mu\nu} \partial_\nu S) T^\alpha \]
    Using the metric inverse property $g_{\mu\alpha} g^{\mu\nu} = \delta_\alpha^\nu$:
    \[ g_{\mu\alpha} l^\mu T^\alpha = \tilde{f}(x) \delta_\alpha^\nu \partial_\nu S T^\alpha = \tilde{f}(x) (\partial_\alpha S) T^\alpha \]
    From Step 1, $T^\alpha \partial_\alpha S = 0$, so the inner product is strictly zero. Thus, $l^\mu$ is orthogonal to all tangent vectors, making it normal to the surface.
\end{enumerate}

\textbf{(b) The Null Property:}
\begin{enumerate}
    \item \textbf{Step 1: Express the norm squared $l^2$.}
    The norm squared of the vector $l^\mu$ is:
    \[ l^2 = g_{\mu\nu} l^\mu l^\nu = g_{\mu\nu} (\tilde{f} g^{\mu\alpha} \partial_\alpha S) (\tilde{f} g^{\nu\beta} \partial_\beta S) \]
    Simplifying using $g_{\mu\nu} g^{\mu\alpha} = \delta_\nu^\alpha$:
    \[ l^2 = \tilde{f}^2 \delta_\nu^\alpha \partial_\alpha S g^{\nu\beta} \partial_\beta S = \tilde{f}^2 g^{\alpha\beta} (\partial_\alpha S) (\partial_\beta S) \]
    
    \item \textbf{Step 2: Evaluate the gradient of $S$.}
    The surface is defined by $S = r - 2M$. Its partial derivatives are:
    \[ \partial_r S = 1, \quad \partial_v S = \partial_\theta S = \partial_\phi S = 0 \]
    \sidenote{We compute these derivatives to find the \textbf{components of the gradient covector} $\partial_\nu S$. Since the surface is defined by the function $S(v, r, \theta, \phi) = r - 2M$, we are measuring how this function varies along each coordinate axis. Because $S$ depends only on $r$, its rate of change along the $v, \theta, \phi$ axes is zero, while its rate of change along the radial axis is exactly 1.}
    Substituting this into the norm squared:
    \[ l^2 = \tilde{f}^2 g^{rr} (\partial_r S)^2 = \tilde{f}^2 g^{rr} \]
    
    \item \textbf{Step 3: Compute the inverse metric component $g^{rr}$.}
    The non-zero components of the metric in the $(v, r)$ sector are $g_{vv} = -(1 - \frac{2M}{r})$ and $g_{vr} = g_{rv} = \frac{1}{2}$. Note that $g_{rr} = 0$.
    The determinant of this $2 \times 2$ block is:
    \[ \det = g_{vv}g_{rr} - (g_{vr})^2 = 0 - \left(\frac{1}{2}\right)^2 = -\frac{1}{4} \]
    The inverse metric component $g^{rr}$ is given by:
    \[ g^{rr} = \frac{g_{vv}}{\det} = \frac{-(1 - \frac{2M}{r})}{-1/4} = 4\left(1 - \frac{2M}{r}\right) \]
    
    \item \textbf{Step 4: Evaluate $l^2$ on the surface $S$.}
    On the hypersurface $S = 0$, we have $r = 2M$.
    Substituting this into $g^{rr}$:
    \[ g^{rr}\Big|_{r=2M} = 4\left(1 - \frac{2M}{2M}\right) = 0 \]
    Therefore, on the surface $S$, the norm squared is $l^2 = \tilde{f}^2 (0) = 0$.
\end{enumerate}

\textbf{Remark on Inverting the Metric:}
A key computational shortcut was used in step (b) to find $g^{rr}$. We only inverted the 2x2 $(v, r)$ sub-matrix. This is permitted because the full 4x4 metric is \textbf{block diagonal}.

The metric $g_{\mu\nu}$ in the $(v, r, \theta, \phi)$ basis has the form:
\[
g_{\mu\nu} = 
\begin{pmatrix}
g_{vv} & g_{vr} & 0 & 0 \\
g_{rv} & g_{rr} & 0 & 0 \\
0 & 0 & g_{\theta\theta} & 0 \\
0 & 0 & 0 & g_{\phi\phi}
\end{pmatrix}
=
\begin{pmatrix}
\mathbf{A} & \mathbf{0} \\
\mathbf{0} & \mathbf{B}
\end{pmatrix}
\]\sidenote{The zeros in the off-diagonal blocks appear because the line element expression $ds^2 = -(1 - \frac{2M}{r}) dv^2 + dv dr + r^2 d\Omega^2$ has no cross-terms mixing the $(v, r)$ coordinates with the angular $(\theta, \phi)$ coordinates (e.g., no $dv\,d\theta$ or $dr\,d\phi$ terms). The components of the metric tensor $g_{\mu\nu}$ are read directly as the coefficients of the $dx^\mu dx^\nu$ terms in the $ds^2$ expression. Similarly, $g_{\theta\phi}=0$ because the angular part $d\Omega^2 = d\theta^2 + \sin^2\theta\,d\phi^2$ contains no $d\theta\,d\phi$ cross-term, making the angular block diagonal as well.}
where $\mathbf{A}$ is the 2x2 $(v, r)$ block and $\mathbf{B}$ is the 2x2 $(\theta, \phi)$ block. The inverse of such a matrix is the matrix of the block inverses:
\[
g^{\mu\nu} = 
\begin{pmatrix}
\mathbf{A}^{-1} & \mathbf{0} \\
\mathbf{0} & \mathbf{B}^{-1}
\end{pmatrix}
\]
Therefore, the components of the inverse metric in the $(v, r)$ sector are determined \textit{only} by the $(v, r)$ block of the original metric. We have:
\[
\mathbf{A} = 
\begin{pmatrix}
g_{vv} & g_{vr} \\
g_{rv} & g_{rr}
\end{pmatrix}
=
\begin{pmatrix}
-(1 - \frac{2M}{r}) & \frac{1}{2} \\
\frac{1}{2} & 0
\end{pmatrix}
\]
The inverse of a generic 2x2 matrix $\begin{pmatrix} a & b \\ c & d \end{pmatrix}$ is $\frac{1}{ad-bc} \begin{pmatrix} d & -b \\ -c & a \end{pmatrix}$. For our matrix $\mathbf{A}$, the determinant is $\det(\mathbf{A}) = (0) \cdot (-(1-\frac{2M}{r})) - (\frac{1}{2})(\frac{1}{2}) = -\frac{1}{4}$.
The inverse is:
\[
\mathbf{A}^{-1} = 
\begin{pmatrix}
g^{vv} & g^{vr} \\
g^{rv} & g^{rr}
\end{pmatrix}
= \frac{1}{-1/4}
\begin{pmatrix}
0 & -\frac{1}{2} \\
-\frac{1}{2} & -(1 - \frac{2M}{r})
\end{pmatrix}
=
\begin{pmatrix}
0 & 2 \\
2 & 4(1 - \frac{2M}{r})
\end{pmatrix}
\]
From this, we can directly read off the component we needed: $g^{rr} = 4(1 - \frac{2M}{r})$. This confirms the result from step (b) and demonstrates why the 4x4 calculation was unnecessary.

\textbf{(c) Killing Symmetry:}
\begin{enumerate}
    \item \textbf{Step 1: The Definition of a Killing Vector Field.}
    A vector field $k$ is a Killing vector field if the metric is invariant under the flow generated by $k$. Mathematically, this is defined by the vanishing of the \textbf{Lie derivative} of the metric tensor with respect to $k$:
    \[ \mathcal{L}_k g_{\mu\nu} = 0 \]
    We can derive the coordinate-independent \textbf{Killing equation} from this:
    \begin{itemize}
        \item Using the formula for the Lie derivative of a $(0,2)$ tensor: $\mathcal{L}_k g_{\mu\nu} = k^\lambda \partial_\lambda g_{\mu\nu} + g_{\lambda\nu} \partial_\mu k^\lambda + g_{\mu\lambda} \partial_\nu k^\lambda$.
        \item Since the metric compatibility condition $\nabla_\lambda g_{\mu\nu} = 0$ holds for the Levi-Civita connection, we can replace partial derivatives with covariant derivatives:
        \[ \mathcal{L}_k g_{\mu\nu} = k^\lambda \nabla_\lambda g_{\mu\nu} + g_{\lambda\nu} \nabla_\mu k^\lambda + g_{\mu\lambda} \nabla_\nu k^\lambda \]
        \item The first term is zero ($k^\lambda \cdot 0 = 0$). Lowering the indices of the remaining terms ($g_{\lambda\nu} \nabla_\mu k^\lambda = \nabla_\mu k_\nu$):
        \[ \mathcal{L}_k g_{\mu\nu} = \nabla_\mu k_\nu + \nabla_\nu k_\mu = 0 \]
    \end{itemize}
    This is the covariant form of the Killing equation. In coordinate-based calculations, we often use the expansion in partial derivatives:
    \[ k^\sigma \partial_\sigma g_{\mu\nu} + g_{\mu\sigma} \partial_\nu k^\sigma + g_{\nu\sigma} \partial_\mu k^\sigma = 0 \]
    This equation must hold for all components $(\mu, \nu)$ for $k$ to be a Killing vector.

    \item \textbf{Step 2: Identify the Components of the Vector Field.}
    We are testing the vector field $k = \frac{\partial}{\partial v}$. In the coordinate basis $(v, r, \theta, \phi)$, this vector has a `1` in the `$v$` position and `0` everywhere else. Its contravariant components are:
    \[ k^v = 1, \quad k^r = 0, \quad k^\theta = 0, \quad k^\phi = 0 \]
    Or more compactly, $k^\mu = \delta^\mu_v$.

    \item \textbf{Step 3: Substitute the Vector Field into the Killing Equation.}
    Since the components of $k^\mu$ are all constants (either 1 or 0), all of their partial derivatives are zero:
    \[ \partial_\nu k^\sigma = 0 \quad \text{for all } \sigma, \nu \]
    Substituting this into the Killing equation from Step 1, the second and third terms vanish:
    \[ k^\sigma \partial_\sigma g_{\mu\nu} + g_{\mu\sigma} (0) + g_{\nu\sigma} (0) = 0 \]
    The equation simplifies dramatically to:
    \[ k^\sigma \partial_\sigma g_{\mu\nu} = 0 \]

    \item \textbf{Step 4: Evaluate the Simplified Equation.}
    Now we substitute the components of $k^\sigma = \delta^\sigma_v$:
    \[ \delta^\sigma_v \partial_\sigma g_{\mu\nu} = \partial_v g_{\mu\nu} = 0 \]
    This gives us the well-known shortcut: for a coordinate basis vector field $\frac{\partial}{\partial x^\alpha}$, it is a Killing vector if and only if the metric components are independent of that coordinate, i.e., $\partial_\alpha g_{\mu\nu} = 0$.

    \item \textbf{Step 5: Verify the Condition for All Metric Components.}
    We must now explicitly check if $\partial_v g_{\mu\nu} = 0$ for all components of the given metric:
    $ds^2 = -(1 - \frac{2M}{r}) dv^2 + dv dr + r^2 (d\theta^2 + \sin^2\theta d\phi^2)$.
    The non-zero metric components are:
    \begin{itemize}
        \item $g_{vv} = -(1 - \frac{2M}{r})$: Depends only on $r$. $\partial_v g_{vv} = 0$.
        \item $g_{vr} = g_{rv} = \frac{1}{2}$: Is a constant. $\partial_v g_{vr} = 0$.
        \item $g_{\theta\theta} = r^2$: Depends only on $r$. $\partial_v g_{\theta\theta} = 0$.
        \item $g_{\phi\phi} = r^2\sin^2\theta$: Depends only on $r$ and $\theta$. $\partial_v g_{\phi\phi} = 0$.
    \end{itemize}
    All other components (e.g., $g_{v\theta}, g_{r\phi}$) are zero, and their partial derivatives with respect to $v$ are also zero.
    
    \item \textbf{Step 6: Conclusion.}
    Since the condition $\partial_v g_{\mu\nu} = 0$ holds for all components $(\mu, \nu)$, the Killing equation is satisfied. Therefore, $\frac{\partial}{\partial v}$ is a Killing vector field. This reflects the physical fact that the spacetime geometry is symmetric under translations in the "advanced time" coordinate $v$.
\end{enumerate}
\end{solution}

\textbf{Remark on Metric Symmetry:}
A fundamental property of the metric tensor used throughout general relativity is its symmetry:
\[ g_{\mu\alpha} = g_{\alpha\mu} \]
This is not a coincidence but a defining characteristic. It stems from the requirement that the inner product (or scalar product) of any two vectors, $u$ and $v$, must be symmetric, i.e., $g(u, v) = g(v, u)$. When written in component form, this directly implies that the components of the metric tensor must be symmetric. As a result, when the metric is represented as a matrix, that matrix is always symmetric about its main diagonal. This reduces the number of independent components in a 4D spacetime from 16 to 10.

\begin{exercise}
\label{physics:2022:5}
The energy momentum tensor for a relativistic quantum field theory is denoted as $\theta^{\mu\nu}$, which is symmetric and conserved.

\begin{enumerate}[label=(\alph*)]
    \item Define new current $s^\mu = x_\nu \theta^{\mu\nu}$ and $K^{\lambda\mu} = x^2 \theta^{\lambda\mu} - 2x^\lambda x_\rho \theta^{\rho\mu}$. Compute $\partial_\mu s^\mu$ and $\partial_\mu K^{\lambda\mu}$, and explain the condition on $\theta^{\mu\nu}$ so that these new currents are conserved.
    \item Consider a scalar field $\sigma(x)$ which transforms under a scale transformation as $\delta \sigma = x^\lambda \partial_\lambda \sigma + f^{-1}$. We have following Lagrangian:
    \begin{equation}
        \mathcal{L} = \mathcal{L}_s - \frac{\mu_0^2}{2} \phi^2 e^{2f\sigma} + \frac{1}{2f^2} \partial_\mu e^{f\sigma} \partial^\mu e^{f\sigma}
    \end{equation}
    The infinitesimal scale transformation on scalar field $\phi$ is $\delta \phi = (1 + x_\lambda \partial^\lambda)\phi$. Here $\mathcal{L}_s$ is scale invariant part of the Lagrangian. Prove that: the above Lagrangian is scale invariant.
    \item Explain why a classically scale invariant Lagrangian for a quantum field theory may fail to be scale invariant quantum mechanically.
\end{enumerate}
\end{exercise}

\begin{solution}
\textbf{Prerequisites / Core Concepts:}
\begin{itemize}
    \item \textbf{Noether's Theorem:} A fundamental theorem stating that every differentiable symmetry of the action of a physical system has a corresponding conservation law. Specifically, a continuous symmetry transformation leads to a \textbf{conserved current}.
    \item \textbf{Conserved Current ($j^\mu$):} A vector field (current) is said to be \textbf{conserved} if its 4-divergence vanishes:
    \[ \partial_\mu j^\mu = \frac{\partial j^0}{\partial t} + \nabla \cdot \mathbf{j} = 0 \]
    Physically, this ensures that the total "charge" $Q = \int j^0 d^3x$ is constant in time ($\frac{dQ}{dt} = 0$), provided the current vanishes at infinity.
    \item \textbf{Energy-Momentum Tensor ($\theta^{\mu\nu}$):} The current corresponding to spacetime translations. It is symmetric ($\theta^{\mu\nu} = \theta^{\nu\mu}$) and conserved ($\partial_\mu \theta^{\mu\nu} = 0$). In this exercise, it serves as the building block for other currents.
    \item \textbf{Trace ($\theta^\mu_\mu$):} The contraction of the energy-momentum tensor with the metric: $\theta^\mu_\mu = g_{\mu\nu} \theta^{\mu\nu}$.
    \item \textbf{Scale Invariance:} A symmetry where the physics of a system remains unchanged under a coordinate rescaling $x^\mu \to \lambda x^\mu$.
    \item \textbf{Minkowski Norm Squared ($x^2$):} The scalar quantity defined as the inner product of the position vector with itself: $x^2 = g_{\mu\nu} x^\mu x^\nu = x_\mu x^\mu = -(x^0)^2 + (x^1)^2 + (x^2)^2 + (x^3)^2$ (in flat spacetime).
    \item \textbf{Conformal Anomaly:} A phenomenon where a classical symmetry (like scale invariance) is broken when the theory is quantized, typically due to the need to introduce a regularization energy scale.
\end{itemize}

\textbf{Intuition / Motivation:}
If a system looks the same regardless of how much you zoom in or out, it has scale invariance. In relativistic field theory, the mathematical condition for this perfect "fractal" behavior is that the trace of the energy-momentum tensor must be zero. We can prove this by constructing currents for scaling ($s^\mu$) and checking when they are conserved. However, the quantum world is messy. To deal with infinite loops in quantum field theory, we must introduce a "ruler" (a momentum scale $\mu$). This ruler breaks the perfect zoom-in symmetry, causing the quantum version of the theory to lose its scale invariance.

\textbf{Logical Step-by-Step Breakdown:}

\textbf{(a) Divergence Calculations for New Currents:}
\begin{enumerate}
    \item \textbf{Step 1: Calculate $\partial_\mu s^\mu$.}
    The scale current is $s^\mu = x_\nu \theta^{\mu\nu}$. Apply the product rule:
    \[ \partial_\mu s^\mu = \partial_\mu (x_\nu \theta^{\mu\nu}) = (\partial_\mu x_\nu) \theta^{\mu\nu} + x_\nu (\partial_\mu \theta^{\mu\nu}) \]
    Since $x_\nu = g_{\nu\alpha} x^\alpha$, we have $\partial_\mu x_\nu = g_{\nu\alpha} \delta^\alpha_\mu = g_{\mu\nu}$.\sidenote{The identity $\partial_\mu x^\alpha = \delta_\mu^\alpha$ reflects the fact that the coordinates $x^0, x^1, x^2, x^3$ are \textbf{independent variables}. In calculus, the partial derivative $\frac{\partial x^i}{\partial x^j}$ is 1 if $i=j$ and 0 if $i \neq j$. This is the definition of the \textbf{Kronecker delta} $\delta_\mu^\alpha$.}
    By energy-momentum conservation, $\partial_\mu \theta^{\mu\nu} = 0$.
    \[ \partial_\mu s^\mu = g_{\mu\nu} \theta^{\mu\nu} = \theta^\mu_\mu \]
    Thus, $s^\mu$ is conserved if and only if the trace $\theta^\mu_\mu$ vanishes.
    
    \item \textbf{Step 2: Calculate $\partial_\mu K^{\lambda\mu}$.}
    The special conformal current is $K^{\lambda\mu} = x^2 \theta^{\lambda\mu} - 2x^\lambda x_\rho \theta^{\rho\mu}$.
    Apply the product rule to the first term:\sidenote{The derivative of the norm squared is $\partial_\mu x^2 = \partial_\mu (g_{\alpha\beta} x^\alpha x^\beta)$. In flat spacetime, the metric is constant, so $\partial_\mu x^2 = g_{\alpha\beta} (\partial_\mu x^\alpha x^\beta + x^\alpha \partial_\mu x^\beta) = g_{\alpha\beta} (\delta_\mu^\alpha x^\beta + x^\alpha \delta_\mu^\beta) = g_{\mu\beta} x^\beta + g_{\alpha\mu} x^\alpha = x_\mu + x_\mu = 2x_\mu$.}
    \[ \partial_\mu (x^2 \theta^{\lambda\mu}) = (\partial_\mu x^2) \theta^{\lambda\mu} + x^2 (\partial_\mu \theta^{\lambda\mu}) = (2x_\mu) \theta^{\lambda\mu} + 0 = 2x_\mu \theta^{\lambda\mu} \]
    Apply the product rule to the second term (with the factor of $-2$):
    \[ \partial_\mu (x^\lambda x_\rho \theta^{\rho\mu}) = (\partial_\mu x^\lambda) x_\rho \theta^{\rho\mu} + x^\lambda (\partial_\mu x_\rho) \theta^{\rho\mu} + x^\lambda x_\rho (\partial_\mu \theta^{\rho\mu}) \]
    Using $\partial_\mu x^\lambda = \delta^\lambda_\mu$, $\partial_\mu x_\rho = g_{\mu\rho}$, and $\partial_\mu \theta^{\rho\mu} = 0$:
    \[ \partial_\mu (x^\lambda x_\rho \theta^{\rho\mu}) = \delta^\lambda_\mu x_\rho \theta^{\rho\mu} + x^\lambda g_{\mu\rho} \theta^{\rho\mu} = x_\rho \theta^{\rho\lambda} + x^\lambda \theta^\mu_\mu \]
    Since $\theta^{\rho\lambda}$ is symmetric, $x_\rho \theta^{\rho\lambda} = x_\mu \theta^{\mu\lambda} = x_\mu \theta^{\lambda\mu}$.
    
    \item \textbf{Step 3: Combine and find the condition.}
    Subtracting the two parts:
    \[ \partial_\mu K^{\lambda\mu} = 2x_\mu \theta^{\lambda\mu} - 2(x_\mu \theta^{\lambda\mu} + x^\lambda \theta^\mu_\mu) = -2x^\lambda \theta^\mu_\mu \]
    Both $\partial_\mu s^\mu = 0$ and $\partial_\mu K^{\lambda\mu} = 0$ are satisfied if and only if the energy-momentum tensor is traceless: $\theta^\mu_\mu = 0$.
\end{enumerate}

\textbf{(b) Classical Scale Invariance of the Lagrangian:}
\begin{enumerate}
    \item \textbf{Step 1: Goal and Methodology.}
    A Lagrangian $\mathcal{L}$ is considered scale-invariant if its corresponding action $S = \int d^4x \mathcal{L}$ is invariant under scale transformations. This is achieved if the Lagrangian density has a canonical scaling dimension of 4. A quantity $Q$ has a scaling dimension $d$ if its variation under the given infinitesimal transformation is $\delta Q = d \cdot Q + x^\lambda \partial_\lambda Q$.\sidenote{This variation is derived from the \textbf{finite scaling law} $\phi'(\lambda x) = \lambda^{-d} \phi(x)$. Expanding to first order in $\lambda = 1+\epsilon$ gives $\phi(x) + \delta \phi(x) + \epsilon x^\lambda \partial_\lambda \phi(x) = (1-d\epsilon)\phi(x)$, which simplifies to $\delta \phi = -\epsilon(d\phi + x^\lambda \partial_\lambda \phi)$. Taking $\epsilon = -1$ yields the standard generator form used here.} Our goal is to show that the full Lagrangian $\mathcal{L}$ satisfies this rule with $d=4$. We will analyze each term separately.

    \item \textbf{Step 2: The Scale-Invariant Part $\mathcal{L}_s$.}
    The problem statement explicitly defines $\mathcal{L}_s$ as the "scale invariant part of the Lagrangian." This means it is given to have a scaling dimension of 4, so we do not need to prove it.
    
    \item \textbf{Step 3: Detailed Analysis of the Mass Term.}
    Let $\mathcal{L}_m = - \frac{\mu_0^2}{2} \phi^2 e^{2f\sigma}$. We calculate its variation $\delta \mathcal{L}_m$ using the product rule, based on the given field variations $\delta\phi = (1+x^\lambda\partial_\lambda)\phi$ and $\delta\sigma = f^{-1} + x^\lambda\partial_\lambda\sigma$.

    First, find the variation of $\phi^2$:
    \[ \delta(\phi^2) = 2\phi\,\delta\phi = 2\phi(1 + x^\lambda\partial_\lambda)\phi = 2\phi^2 + 2\phi x^\lambda\partial_\lambda\phi = 2\phi^2 + x^\lambda\partial_\lambda(\phi^2) \]
    This shows $\phi^2$ has a scaling dimension of 2.

    Next, find the variation of $e^{2f\sigma}$:
    \[ \delta(e^{2f\sigma}) = \frac{\partial(e^{2f\sigma})}{\partial\sigma}\delta\sigma = (2f e^{2f\sigma}) \delta\sigma = 2f e^{2f\sigma} (f^{-1} + x^\lambda\partial_\lambda\sigma) = 2e^{2f\sigma} + 2f e^{2f\sigma} x^\lambda\partial_\lambda\sigma \]
    Recognizing that $x^\lambda \partial_\lambda(e^{2f\sigma}) = 2f e^{2f\sigma} x^\lambda\partial_\lambda\sigma$, we get:
    \[ \delta(e^{2f\sigma}) = 2e^{2f\sigma} + x^\lambda\partial_\lambda(e^{2f\sigma}) \]
    This shows $e^{2f\sigma}$ also has a scaling dimension of 2.

    Now, combine these using the product rule for variations:
    \[ \delta \mathcal{L}_m = - \frac{\mu_0^2}{2} \left[ \delta(\phi^2) e^{2f\sigma} + \phi^2 \delta(e^{2f\sigma}) \right] \]
    \[ \delta \mathcal{L}_m = - \frac{\mu_0^2}{2} \left[ (2\phi^2 + x^\lambda\partial_\lambda(\phi^2)) e^{2f\sigma} + \phi^2 (2e^{2f\sigma} + x^\lambda\partial_\lambda(e^{2f\sigma})) \right] \]
    \[ \delta \mathcal{L}_m = - \frac{\mu_0^2}{2} \left[ 4\phi^2 e^{2f\sigma} + \left( (x^\lambda\partial_\lambda(\phi^2)) e^{2f\sigma} + \phi^2 x^\lambda\partial_\lambda(e^{2f\sigma}) \right) \right] \]
    Using the product rule for derivatives on the last two terms:
    \[ \delta \mathcal{L}_m = - \frac{\mu_0^2}{2} \left[ 4\phi^2 e^{2f\sigma} + x^\lambda\partial_\lambda(\phi^2 e^{2f\sigma}) \right] \]
    Distributing the constant factor gives:
    \[ \delta \mathcal{L}_m = 4 \left( - \frac{\mu_0^2}{2} \phi^2 e^{2f\sigma} \right) + x^\lambda\partial_\lambda \left( - \frac{\mu_0^2}{2} \phi^2 e^{2f\sigma} \right) = 4\mathcal{L}_m + x^\lambda\partial_\lambda \mathcal{L}_m \]
    This confirms that $\mathcal{L}_m$ has a scaling dimension of 4.

    \item \textbf{Step 4: Detailed Analysis of the Kinetic Term.}
    Let $\mathcal{L}_\sigma = \frac{1}{2f^2} \partial_\mu e^{f\sigma} \partial^\mu e^{f\sigma}$. Let's define an auxiliary field $\Psi = e^{f\sigma}$.
    First, let's find the scaling dimension of $\Psi$:
    \[ \delta\Psi = \frac{\partial(e^{f\sigma})}{\partial\sigma}\delta\sigma = f e^{f\sigma}\delta\sigma = f\Psi(f^{-1} + x^\lambda\partial_\lambda\sigma) = \Psi + f\Psi x^\lambda\partial_\lambda\sigma = \Psi + x^\lambda\partial_\lambda\Psi \]
    This shows $\Psi$ has a scaling dimension of 1, just like $\phi$.

    Next, we use the standard transformation rule for the derivative of a field $A$ with dimension $d_A$: its derivative $\partial_\mu A$ transforms with dimension $d_A+1$, so $\delta(\partial_\mu A) = (d_A+1)\partial_\mu A + x^\lambda\partial_\lambda(\partial_\mu A)$.
    Since $\Psi$ has dimension 1, its derivative $\partial_\mu\Psi$ has dimension $1+1=2$. Thus:
    \[ \delta(\partial_\mu\Psi) = 2(\partial_\mu\Psi) + x^\lambda\partial_\lambda(\partial_\mu\Psi) \]
    
    Now we can compute the variation of $\mathcal{L}_\sigma = \frac{1}{2f^2}g^{\mu\nu}(\partial_\mu\Psi)(\partial_\nu\Psi)$:
    \[ \delta\mathcal{L}_\sigma = \frac{1}{2f^2}g^{\mu\nu} \left[ \delta(\partial_\mu\Psi)(\partial_\nu\Psi) + (\partial_\mu\Psi)\delta(\partial_\nu\Psi) \right] \]
    Since the expression is symmetric in $\mu, \nu$, this becomes:
    \[ \delta\mathcal{L}_\sigma = \frac{1}{f^2}g^{\mu\nu} (\partial_\nu\Psi) \delta(\partial_\mu\Psi) \]
    Substitute the transformation rule for $\delta(\partial_\mu\Psi)$:
    \[ \delta\mathcal{L}_\sigma = \frac{1}{f^2}g^{\mu\nu} (\partial_\nu\Psi) \left[ 2(\partial_\mu\Psi) + x^\lambda\partial_\lambda(\partial_\mu\Psi) \right] \]
    \[ \delta\mathcal{L}_\sigma = \frac{2}{f^2}g^{\mu\nu}(\partial_\nu\Psi)(\partial_\mu\Psi) + \frac{1}{f^2}g^{\mu\nu} (\partial_\nu\Psi) x^\lambda\partial_\lambda(\partial_\mu\Psi) \]
    The first term is exactly $4\mathcal{L}_\sigma$. The second term is $x^\lambda\partial_\lambda\mathcal{L}_\sigma$ by the product rule for derivatives:
    \[ x^\lambda\partial_\lambda\mathcal{L}_\sigma = x^\lambda\partial_\lambda \left( \frac{1}{2f^2}g^{\mu\nu}(\partial_\mu\Psi)(\partial_\nu\Psi) \right) = \frac{1}{2f^2}g^{\mu\nu}x^\lambda \left[ (\partial_\lambda\partial_\mu\Psi)(\partial_\nu\Psi) + (\partial_\mu\Psi)(\partial_\lambda\partial_\nu\Psi) \right] \]
    \[ x^\lambda\partial_\lambda\mathcal{L}_\sigma = \frac{1}{f^2}g^{\mu\nu}x^\lambda (\partial_\nu\Psi) (\partial_\lambda\partial_\mu\Psi) \]
    Therefore, we have shown:
    \[ \delta\mathcal{L}_\sigma = 4\mathcal{L}_\sigma + x^\lambda\partial_\lambda\mathcal{L}_\sigma \]
    This confirms that $\mathcal{L}_\sigma$ also has a scaling dimension of 4.

    \item \textbf{Step 5: Conclusion.}
    Since $\mathcal{L}_s$, $\mathcal{L}_m$, and $\mathcal{L}_\sigma$ all have a scaling dimension of 4, their sum, the total Lagrangian $\mathcal{L}$, also has a scaling dimension of 4. A Lagrangian with scaling dimension 4 yields a scale-invariant action, completing the proof.
\end{enumerate}

\textbf{(c) The Quantum Scale Anomaly:}
\begin{enumerate}
    \item \textbf{Step 1: The necessity of regularization.}
    In quantum field theory, loop diagrams produce divergent integrals. To make sense of these infinities, we must regulate the theory (e.g., via dimensional regularization or a momentum cutoff).
    
    \item \textbf{Step 2: Introduction of a mass scale.}
    Every regularization scheme requires the introduction of an arbitrary mass/momentum scale $\mu$ (the renormalization scale) to keep the units consistent or to set the cutoff limit.
    
    \item \textbf{Step 3: Breaking of scale invariance.}
    Because the coupling constants now depend on this introduced scale $\mu$ (running couplings), the theory is no longer invariant under arbitrary rescaling of coordinates $x \to \lambda x$. This translates to a non-zero trace of the energy momentum tensor at the quantum level ($\langle \theta^\mu_\mu \rangle \neq 0$), a phenomenon known as the trace anomaly or conformal anomaly.
\end{enumerate}

\begin{remark}[QFT Regularization Techniques]
To compute divergent loop integrals in quantum field theory, we use \textbf{regularization} to make the results finite before removing the regulator in the final step. Two primary methods mentioned in the proof are:
\begin{itemize}
    \item \textbf{Momentum Cutoff:} A direct approach where we restrict the range of integration in momentum space to $|k| < \Lambda$. The parameter $\Lambda$ acts as a "maximum possible momentum" (minimum possible distance). While intuitive, this method often breaks symmetries like gauge invariance and Lorentz invariance.
    \item \textbf{Dimensional Regularization (DimReg):} A more sophisticated method that generalizes the spacetime dimension from $d=4$ to $d = 4 - \epsilon$, where $\epsilon$ is a small parameter. The divergences appear as poles (terms like $1/\epsilon$) as $\epsilon \to 0$. Because it does not involve a hard cutoff in momentum, it preserves most fundamental symmetries of the theory.
\end{itemize}
Crucially, both methods introduce an artificial energy scale (either $\Lambda$ or the mass scale $\mu$ required for dimensional consistency in DimReg) that breaks the classical scale invariance of the Lagrangian.
\end{remark}
\end{solution}

\begin{exercise}
\label{physics:2022:6}
Consider following Lagrangian for $N$ scalar fields $\phi^a, a = 1, \dots, N$:
\begin{equation}
    \mathcal{L} = \frac{1}{2} \partial_\mu \phi^a \partial^\mu \phi^a - \frac{1}{2} \mu_0^2 \phi^a \phi^a - \frac{1}{8} \lambda_0 (\phi^a \phi^a)^2
\end{equation}
Here the repeated index implies the summation over the index.

\begin{enumerate}[label=(\alph*)]
    \item Write down the propagator and interaction vertex for this model, and write down four point Feynman diagrams up to one loop level.
    \item Define $g_0 = \lambda_0 N$, and compute the order in $g_0$ and $N$ for all the diagrams listed in last question. If we fix the coupling $g_0$, and let $N$ go to infinity, list the leading order Feynman diagrams in $1/N$.
\end{enumerate}
\end{exercise}

\begin{solution}
\textbf{Prerequisites / Core Concepts:}
\begin{itemize}
    \item \textbf{Feynman Rules:} A set of graphical rules for calculating scattering amplitudes in quantum field theory.
    \begin{itemize}
        \item \textbf{Propagator:} The propagator $D(x-y)$ is the vacuum expectation value of the time-ordered product of two fields, $\langle 0 | T\{\phi(x)\phi(y)\} | 0 \rangle$. It represents the amplitude for a particle to propagate between spacetime points $x$ and $y$. In momentum space, for a scalar field with mass $\mu_0$, it is the inverse of the free-field operator:
        \[ D_{ab}(p) = \frac{i \delta_{ab}}{p^2 - \mu_0^2 + i\epsilon} \]
        The $i\epsilon$ term specifies the correct integration contour for the Green's function.
        \item \textbf{Interaction Vertex:} A vertex represents a fundamental interaction. Its mathematical value is derived from the interaction term in the Lagrangian, $\mathcal{L}_{\text{int}}$. For a 4-point vertex from a $\lambda_0\phi^4$ interaction, the rule is given by the functional derivative of the interaction action, multiplied by $i$:
        \[ \text{Vertex} \propto i \frac{\delta^4}{\delta\phi(x_1)\dots\delta\phi(x_4)} \int d^4x \mathcal{L}_{\text{int}} \]
        For the Lagrangian term $-\frac{\lambda_0}{8}(\phi^a\phi^a)^2$, this yields the momentum-space rule (a tensor for the indices $a,b,c,d$):
        \[ V_{abcd} = -i \lambda_0 (\delta_{ab}\delta_{cd} + \delta_{ac}\delta_{bd} + \delta_{ad}\delta_{bc}) \]
    \end{itemize}
    \item \textbf{Large $N$ Expansion ('t Hooft Limit):} A perturbative expansion technique for theories with an $O(N)$ or $U(N)$ symmetry group, where the expansion parameter is $1/N$. The limit is formally defined as:
    \[ N \to \infty, \quad \lambda_0 \to 0, \quad \text{with the 't Hooft coupling } g_0 = \lambda_0 N \text{ held fixed.} \]
    The amplitude for any Feynman diagram can then be classified by its power of $N$. An amplitude $\mathcal{A}$ for a given process is expressed as a series:
    \[ \mathcal{A}(g_0, N) = \sum_{k=0}^\infty \frac{1}{N^k} f_k(g_0) \]
    The goal is to find the leading-order diagrams, which are those that survive in the $N\to\infty$ limit (or have the lowest power of $1/N$). These are often a tractable, infinite subclass of diagrams (e.g., "planar" diagrams).
\end{itemize}

\textbf{Intuition / Motivation:}
Imagine calculating the probability of a rumor spreading in a city with $N$ different professions. If a rumor can pass between any two people of the same profession, the most dominant way it spreads involves chains of people passing it along in closed circles. Even if the individual chance of sharing (the coupling $\lambda_0$) is small, the sheer number of professions ($N$) makes these "closed loop" chains dominate over everything else. In quantum field theory, this is the essence of the large-$N$ limit: diagrams that maximize the number of independent "index loops" are the most important.

\textbf{Logical Step-by-Step Breakdown:}

\textbf{(a) Feynman Rules and Diagrams up to One Loop:}
\begin{enumerate}
    \item \textbf{Step 1: Determine the propagator.}
    The free part of the Lagrangian is $\frac{1}{2} \partial_\mu \phi^a \partial^\mu \phi^a - \frac{1}{2} \mu_0^2 \phi^a \phi^a$. This is simply $N$ independent scalar fields. The propagator for a field of species $a$ to species $b$ requires them to be the same species, so it carries a Kronecker delta $\delta_{ab}$:
    \[ D_{ab}(p) = \frac{i \delta_{ab}}{p^2 - \mu_0^2 + i\epsilon} \]
    
    \item \textbf{Step 2: Determine the interaction vertex.}
    The interaction term is $-\frac{1}{8} \lambda_0 (\phi^a \phi^a)(\phi^b \phi^b)$.
    Taking four functional derivatives to find the 4-point vertex rule (and multiplying by $i$), we must account for all ways to pair the four incoming fields. The factor of $1/8$ is cancelled by the $8$ permutations of pairing indices. The vertex rule for incoming fields with indices $a, b, c, d$ is:
    \[ V_{abcd} = -i \lambda_0 (\delta_{ab}\delta_{cd} + \delta_{ac}\delta_{bd} + \delta_{ad}\delta_{bc}) \]
    
    \item \textbf{Step 3: List the four-point diagrams.}
    \begin{itemize}
        \item \textbf{Tree Level:} There is exactly 1 diagram consisting of a single vertex connecting four external lines.
        \item \textbf{One-Loop Level:} There are 3 diagrams, corresponding to the $s$, $t$, and $u$ scattering channels. Each consists of two vertices connected by two internal propagators forming a "bubble".
    \end{itemize}
\end{enumerate}

\textbf{(b) Large $N$ Scaling and Leading Diagrams:}
\begin{enumerate}
    \item \textbf{Step 1: Set up the 't Hooft coupling.}
    We define $g_0 = \lambda_0 N$. To hold $g_0$ fixed as $N \to \infty$, we must have the bare coupling scale as $\lambda_0 = \frac{g_0}{N}$.
    
    \item \textbf{Step 2: Scale of the tree-level diagram.}
    The tree-level diagram has exactly one vertex and no internal loops.
    Its scaling is simply proportional to the coupling:
    \[ \text{Tree Diagram} \sim \lambda_0 = \frac{g_0}{N} \sim \mathcal{O}\left(\frac{1}{N}\right) \]
    
    \item \textbf{Step 3: Scale of the one-loop bubble diagrams.}
    A one-loop diagram has two vertices, so it carries a factor of $\lambda_0^2 = \frac{g_0^2}{N^2}$.
    However, we must also count the number of free internal indices. 
    In the $s$-channel diagram (where initial states $a, b$ annihilate into the loop and recreate $c, d$), if we use the $\delta_{ab}\delta_{ii}$ term from the first vertex and $\delta_{cd}\delta_{ii}$ from the second, the internal propagators both carry the index $i$.
    Summing over this free internal index gives $\sum_{i=1}^N 1 = N$.
    Thus, this diagram scales as:
    \[ \text{One-loop } s\text{-channel} \sim \lambda_0^2 \times N = \frac{g_0^2}{N^2} \times N = \frac{g_0^2}{N} \sim \mathcal{O}\left(\frac{1}{N}\right) \]
    Note that $t$ and $u$ channel diagrams do not have a free internal index sum that yields $N$ (their internal indices are constrained by the external ones), so they scale as $\lambda_0^2 \sim \mathcal{O}(1/N^2)$ and are suppressed.
    
    \item \textbf{Step 4: Generalize to leading order diagrams.}
    Every time we add a new "bubble" to form a chain in the $s$-channel, we add one more vertex ($\sim 1/N$) and one more closed index loop ($\sim N$). The factors of $N$ cancel out!
    For a chain with $L$ loops, there are $L+1$ vertices and $L$ closed index sums:
    \[ \text{Bubble chain} \sim (\lambda_0)^{L+1} N^L = \left(\frac{g_0}{N}\right)^{L+1} N^L = \frac{g_0^{L+1}}{N} \sim \mathcal{O}\left(\frac{1}{N}\right) \]
    Therefore, if we fix $g_0$ and let $N \to \infty$, the leading order diagrams (all contributing at $\mathcal{O}(1/N)$) are the tree-level diagram and the infinite series of $s$-channel "bubble chain" diagrams.
\end{enumerate}
\end{solution}
